% Ngân hàng đề Toán 12 (KNTT) -- Chương II, Bài 6, Dạng 03
% Dạng: Tích vô hướng và góc giữa hai vectơ trong không gian
%% Phương án ĐẦU TIÊN trong \cauTN luôn là đáp án đúng (gói trondethi tự đảo).

\cauTN[id=C2-B6-D03-NB-001,muc=NB,loigiai={Biểu diễn hai vectơ theo ba cạnh xuất phát từ một đỉnh: $\overrightarrow{DG}=\vec{a}+\vec{c}$, $\overrightarrow{FE}=-\vec{a}$, rồi dùng $\vec{x}\cdot\vec{y}$ theo các tích vô hướng cơ bản ta được $-a^2$.}]{Cho hình lập phương $ABCD.EFGH$ có cạnh bằng $a$. Tích vô hướng $\overrightarrow{DG}\cdot \overrightarrow{FE}$ bằng}
  {$-a^2$.}
  {$0$.}{$-2a^2$.}{$a^2$.}

\cauTN[id=C2-B6-D03-NB-002,muc=NB,loigiai={Biểu diễn hai vectơ theo ba cạnh xuất phát từ một đỉnh: $\overrightarrow{FE'}=-\vec{a}+\vec{c}$, $\overrightarrow{G'E}=-\vec{a}-\vec{b}-\vec{c}$, rồi dùng $\vec{x}\cdot\vec{y}$ theo các tích vô hướng cơ bản ta được $0$.}]{Cho hình lập phương $EFGH.E'F'G'H'$ có cạnh bằng $2$. Tích vô hướng $\overrightarrow{FE'}\cdot \overrightarrow{G'E}$ bằng}
  {$0$.}
  {$1$.}{$-1$.}{$2$.}

\cauTN[id=C2-B6-D03-NB-003,muc=NB,loigiai={Biểu diễn hai vectơ theo ba cạnh xuất phát từ một đỉnh: $\overrightarrow{MQ}=\vec{c}$, $\overrightarrow{QP}=\vec{b}-\vec{c}$, rồi dùng $\vec{x}\cdot\vec{y}$ theo các tích vô hướng cơ bản ta được $-2$.}]{Cho tứ diện đều $MNPQ$ có cạnh bằng $2$. Tích vô hướng $\overrightarrow{MQ}\cdot \overrightarrow{QP}$ bằng}
  {$-2$.}
  {$-1$.}{$-3$.}{$2$.}

\cauTN[id=C2-B6-D03-NB-004,muc=NB,loigiai={Biểu diễn hai vectơ theo ba cạnh xuất phát từ một đỉnh: $\overrightarrow{Q'Q}=-\vec{c}$, $\overrightarrow{MN'}=\vec{a}+\vec{c}$, rồi dùng $\vec{x}\cdot\vec{y}$ theo các tích vô hướng cơ bản ta được $-4$.}]{Cho hình lập phương $MNPQ.M'N'P'Q'$ có cạnh bằng $2$. Tích vô hướng $\overrightarrow{Q'Q}\cdot \overrightarrow{MN'}$ bằng}
  {$-4$.}
  {$-3$.}{$-5$.}{$4$.}

\cauTN[id=C2-B6-D03-NB-005,muc=NB,loigiai={Biểu diễn hai vectơ theo ba cạnh xuất phát từ một đỉnh: $\overrightarrow{PN}=\vec{a}-\vec{b}$, $\overrightarrow{MQ}=\vec{c}$, rồi dùng $\vec{x}\cdot\vec{y}$ theo các tích vô hướng cơ bản ta được $0$.}]{Cho tứ diện đều $MNPQ$ có cạnh bằng $1$. Tích vô hướng $\overrightarrow{PN}\cdot \overrightarrow{MQ}$ bằng}
  {$0$.}
  {$1$.}{$-1$.}{$2$.}

\cauTN[id=C2-B6-D03-NB-006,muc=NB,loigiai={Biểu diễn hai vectơ theo ba cạnh xuất phát từ một đỉnh: $\overrightarrow{BC}=-\vec{a}+\vec{b}$, $\overrightarrow{CD}=-\vec{b}+\vec{c}$, rồi dùng $\vec{x}\cdot\vec{y}$ theo các tích vô hướng cơ bản ta được $-\dfrac{9}{2}$.}]{Cho tứ diện đều $ABCD$ có cạnh bằng $3$. Tích vô hướng $\overrightarrow{BC}\cdot \overrightarrow{CD}$ bằng}
  {$-\dfrac{9}{2}$.}
  {$-\dfrac{7}{2}$.}{$-\dfrac{11}{2}$.}{$\dfrac{9}{2}$.}

\cauTN[id=C2-B6-D03-NB-007,muc=NB,loigiai={Biểu diễn hai vectơ theo ba cạnh xuất phát từ một đỉnh: $\overrightarrow{DB}=\vec{a}-\vec{c}$, $\overrightarrow{AB}=\vec{a}$, rồi dùng $\vec{x}\cdot\vec{y}$ theo các tích vô hướng cơ bản ta được $\dfrac{9}{2}$.}]{Cho tứ diện đều $ABCD$ có cạnh bằng $3$. Tích vô hướng $\overrightarrow{DB}\cdot \overrightarrow{AB}$ bằng}
  {$\dfrac{9}{2}$.}
  {$\dfrac{11}{2}$.}{$\dfrac{7}{2}$.}{$-\dfrac{9}{2}$.}

\cauTN[id=C2-B6-D03-NB-008,muc=NB,loigiai={Biểu diễn hai vectơ theo ba cạnh xuất phát từ một đỉnh: $\overrightarrow{CS}=-\vec{c}$, $\overrightarrow{AC}=-\vec{a}+\vec{c}$, rồi dùng $\vec{x}\cdot\vec{y}$ theo các tích vô hướng cơ bản ta được $-\dfrac{9}{2}$.}]{Cho tứ diện đều $SABC$ có cạnh bằng $3$. Tích vô hướng $\overrightarrow{CS}\cdot \overrightarrow{AC}$ bằng}
  {$-\dfrac{9}{2}$.}
  {$-\dfrac{7}{2}$.}{$-\dfrac{11}{2}$.}{$\dfrac{9}{2}$.}

\cauTN[id=C2-B6-D03-NB-009,muc=NB,loigiai={Biểu diễn hai vectơ theo ba cạnh xuất phát từ một đỉnh: $\overrightarrow{HF}=\vec{a}-\vec{b}$, $\overrightarrow{CB}=-\vec{b}$, rồi dùng $\vec{x}\cdot\vec{y}$ theo các tích vô hướng cơ bản ta được $a^2$.}]{Cho hình lập phương $ABCD.EFGH$ có cạnh bằng $a$. Tích vô hướng $\overrightarrow{HF}\cdot \overrightarrow{CB}$ bằng}
  {$a^2$.}
  {$2a^2$.}{$0$.}{$-a^2$.}

\cauTN[id=C2-B6-D03-NB-010,muc=NB,loigiai={Biểu diễn hai vectơ theo ba cạnh xuất phát từ một đỉnh: $\overrightarrow{A'B'}=\vec{a}$, $\overrightarrow{AB'}=\vec{a}+\vec{c}$, rồi dùng $\vec{x}\cdot\vec{y}$ theo các tích vô hướng cơ bản ta được $1$.}]{Cho hình lập phương $ABCD.A'B'C'D'$ có cạnh bằng $1$. Tích vô hướng $\overrightarrow{A'B'}\cdot \overrightarrow{AB'}$ bằng}
  {$1$.}
  {$2$.}{$0$.}{$-1$.}

\cauTN[id=C2-B6-D03-NB-011,muc=NB,loigiai={Biểu diễn hai vectơ theo ba cạnh xuất phát từ một đỉnh: $\overrightarrow{SA}=\vec{a}$, $\overrightarrow{SC}=\vec{c}$, rồi dùng $\vec{x}\cdot\vec{y}$ theo các tích vô hướng cơ bản ta được $\dfrac{9}{2}$.}]{Cho tứ diện đều $SABC$ có cạnh bằng $3$. Tích vô hướng $\overrightarrow{SA}\cdot \overrightarrow{SC}$ bằng}
  {$\dfrac{9}{2}$.}
  {$\dfrac{11}{2}$.}{$\dfrac{7}{2}$.}{$-\dfrac{9}{2}$.}

\cauTN[id=C2-B6-D03-NB-012,muc=NB,loigiai={Biểu diễn hai vectơ theo ba cạnh xuất phát từ một đỉnh: $\overrightarrow{PQ}=-\vec{b}+\vec{c}$, $\overrightarrow{QN}=\vec{a}-\vec{c}$, rồi dùng $\vec{x}\cdot\vec{y}$ theo các tích vô hướng cơ bản ta được $-2$.}]{Cho tứ diện đều $MNPQ$ có cạnh bằng $2$. Tích vô hướng $\overrightarrow{PQ}\cdot \overrightarrow{QN}$ bằng}
  {$-2$.}
  {$-1$.}{$-3$.}{$2$.}

\cauTN[id=C2-B6-D03-NB-013,muc=NB,loigiai={Ta có $\overrightarrow{HE}\cdot \overrightarrow{AF}=0$, $|\overrightarrow{HE}|^2=9$, $|\overrightarrow{AF}|^2=17$ nên $\cos=\dfrac{0}{\sqrt{9\cdot 17}}$, suy ra góc bằng $90^\circ$.}]{Cho hình hộp chữ nhật $ABCD.EFGH$ có $AB=4$, $AD=3$, $AE=1$. Góc giữa hai vectơ $\overrightarrow{HE}$ và $\overrightarrow{AF}$ bằng}
  {$90^\circ$.}
  {$135^\circ$.}{$150^\circ$.}{$120^\circ$.}

\cauTN[id=C2-B6-D03-NB-014,muc=NB,loigiai={Ta có $\overrightarrow{QQ'}\cdot \overrightarrow{P'M'}=0$, $|\overrightarrow{QQ'}|^2=9$, $|\overrightarrow{P'M'}|^2=18$ nên $\cos=\dfrac{0}{\sqrt{9\cdot 18}}$, suy ra góc bằng $90^\circ$.}]{Cho hình lập phương $MNPQ.M'N'P'Q'$ có cạnh bằng $3$. Góc giữa hai vectơ $\overrightarrow{QQ'}$ và $\overrightarrow{P'M'}$ bằng}
  {$90^\circ$.}
  {$45^\circ$.}{$120^\circ$.}{$150^\circ$.}

\cauTN[id=C2-B6-D03-NB-015,muc=NB,loigiai={Ta có $\overrightarrow{DD'}\cdot \overrightarrow{A'C'}=0$, $|\overrightarrow{DD'}|^2=4$, $|\overrightarrow{A'C'}|^2=2$ nên $\cos=\dfrac{0}{\sqrt{4\cdot 2}}$, suy ra góc bằng $90^\circ$.}]{Cho hình hộp chữ nhật $ABCD.A'B'C'D'$ có $AB=1$, $AD=1$, $AA'=2$. Góc giữa hai vectơ $\overrightarrow{DD'}$ và $\overrightarrow{A'C'}$ bằng}
  {$90^\circ$.}
  {$30^\circ$.}{$45^\circ$.}{$120^\circ$.}

\cauTN[id=C2-B6-D03-NB-016,muc=NB,loigiai={Ta có $\overrightarrow{CB}\cdot \overrightarrow{AS}=0$, $|\overrightarrow{CB}|^2=a^2$, $|\overrightarrow{AS}|^2=a^2$ nên $\cos=\dfrac{0}{\sqrt{a^2\cdot a^2}}$, suy ra góc bằng $90^\circ$.}]{Cho tứ diện đều $SABC$ có cạnh bằng $a$. Góc giữa hai vectơ $\overrightarrow{CB}$ và $\overrightarrow{AS}$ bằng}
  {$90^\circ$.}
  {$30^\circ$.}{$45^\circ$.}{$120^\circ$.}

\cauTN[id=C2-B6-D03-NB-017,muc=NB,loigiai={Ta có $\overrightarrow{PP'}\cdot \overrightarrow{N'P}=-1$, $|\overrightarrow{PP'}|^2=1$, $|\overrightarrow{N'P}|^2=2$ nên $\cos=\dfrac{-1}{\sqrt{1\cdot 2}}$, suy ra góc bằng $135^\circ$.}]{Cho hình lập phương $MNPQ.M'N'P'Q'$ có cạnh bằng $1$. Góc giữa hai vectơ $\overrightarrow{PP'}$ và $\overrightarrow{N'P}$ bằng}
  {$135^\circ$.}
  {$150^\circ$.}{$45^\circ$.}{$30^\circ$.}

\cauTN[id=C2-B6-D03-NB-018,muc=NB,loigiai={Ta có $\overrightarrow{HG}\cdot \overrightarrow{FH}=-1$, $|\overrightarrow{HG}|^2=1$, $|\overrightarrow{FH}|^2=2$ nên $\cos=\dfrac{-1}{\sqrt{1\cdot 2}}$, suy ra góc bằng $135^\circ$.}]{Cho hình lập phương $EFGH.E'F'G'H'$ có cạnh bằng $1$. Góc giữa hai vectơ $\overrightarrow{HG}$ và $\overrightarrow{FH}$ bằng}
  {$135^\circ$.}
  {$30^\circ$.}{$45^\circ$.}{$150^\circ$.}

\cauTN[id=C2-B6-D03-NB-019,muc=NB,loigiai={Ta có $\overrightarrow{N'P'}\cdot \overrightarrow{N'M}=0$, $|\overrightarrow{N'P'}|^2=9$, $|\overrightarrow{N'M}|^2=8$ nên $\cos=\dfrac{0}{\sqrt{9\cdot 8}}$, suy ra góc bằng $90^\circ$.}]{Cho hình hộp chữ nhật $MNPQ.M'N'P'Q'$ có $MN=2$, $MQ=3$, $MM'=2$. Góc giữa hai vectơ $\overrightarrow{N'P'}$ và $\overrightarrow{N'M}$ bằng}
  {$90^\circ$.}
  {$120^\circ$.}{$60^\circ$.}{$30^\circ$.}

\cauTN[id=C2-B6-D03-NB-020,muc=NB,loigiai={Ta có $\overrightarrow{AE}\cdot \overrightarrow{FE}=0$, $|\overrightarrow{AE}|^2=1$, $|\overrightarrow{FE}|^2=1$ nên $\cos=\dfrac{0}{\sqrt{1\cdot 1}}$, suy ra góc bằng $90^\circ$.}]{Cho hình lập phương $ABCD.EFGH$ có cạnh bằng $1$. Góc giữa hai vectơ $\overrightarrow{AE}$ và $\overrightarrow{FE}$ bằng}
  {$90^\circ$.}
  {$120^\circ$.}{$45^\circ$.}{$135^\circ$.}

\cauTN[id=C2-B6-D03-NB-021,muc=NB,loigiai={Ta có $\overrightarrow{E'E}\cdot \overrightarrow{GH'}=-a^2$, $|\overrightarrow{E'E}|^2=a^2$, $|\overrightarrow{GH'}|^2=2a^2$ nên $\cos=\dfrac{-a^2}{\sqrt{a^2\cdot 2a^2}}$, suy ra góc bằng $135^\circ$.}]{Cho hình lập phương $EFGH.E'F'G'H'$ có cạnh bằng $a$. Góc giữa hai vectơ $\overrightarrow{E'E}$ và $\overrightarrow{GH'}$ bằng}
  {$135^\circ$.}
  {$90^\circ$.}{$30^\circ$.}{$60^\circ$.}

\cauTN[id=C2-B6-D03-NB-022,muc=NB,loigiai={Ta có $\overrightarrow{NQ}\cdot \overrightarrow{PM}=0$, $|\overrightarrow{NQ}|^2=4$, $|\overrightarrow{PM}|^2=4$ nên $\cos=\dfrac{0}{\sqrt{4\cdot 4}}$, suy ra góc bằng $90^\circ$.}]{Cho tứ diện đều $MNPQ$ có cạnh bằng $2$. Góc giữa hai vectơ $\overrightarrow{NQ}$ và $\overrightarrow{PM}$ bằng}
  {$90^\circ$.}
  {$150^\circ$.}{$60^\circ$.}{$30^\circ$.}

\cauTN[id=C2-B6-D03-NB-023,muc=NB,loigiai={Ta có $\overrightarrow{EF}\cdot \overrightarrow{GF}=0$, $|\overrightarrow{EF}|^2=1$, $|\overrightarrow{GF}|^2=1$ nên $\cos=\dfrac{0}{\sqrt{1\cdot 1}}$, suy ra góc bằng $90^\circ$.}]{Cho hình lập phương $ABCD.EFGH$ có cạnh bằng $1$. Góc giữa hai vectơ $\overrightarrow{EF}$ và $\overrightarrow{GF}$ bằng}
  {$90^\circ$.}
  {$135^\circ$.}{$30^\circ$.}{$60^\circ$.}

\cauTN[id=C2-B6-D03-NB-024,muc=NB,loigiai={Ta có $\overrightarrow{AB}\cdot \overrightarrow{CD}=0$, $|\overrightarrow{AB}|^2=9$, $|\overrightarrow{CD}|^2=9$ nên $\cos=\dfrac{0}{\sqrt{9\cdot 9}}$, suy ra góc bằng $90^\circ$.}]{Cho tứ diện đều $ABCD$ có cạnh bằng $3$. Góc giữa hai vectơ $\overrightarrow{AB}$ và $\overrightarrow{CD}$ bằng}
  {$90^\circ$.}
  {$30^\circ$.}{$150^\circ$.}{$60^\circ$.}

\cauTN[id=C2-B6-D03-NB-025,muc=NB,loigiai={Hai vectơ vuông góc khi tích vô hướng bằng $0$. Với $\overrightarrow{DB}$ và $\overrightarrow{CA}$ ta tính được tích vô hướng bằng $0$; các cặp còn lại có tích vô hướng khác $0$.}]{Cho tứ diện đều $ABCD$ có cạnh bằng $1$. Cặp vectơ nào dưới đây vuông góc với nhau?}
  {$\overrightarrow{DB}$ và $\overrightarrow{CA}$.}
  {$\overrightarrow{DB}$ và $\overrightarrow{BA}$.}{$\overrightarrow{AB}$ và $\overrightarrow{BC}$.}{$\overrightarrow{AD}$ và $\overrightarrow{BA}$.}

\cauTN[id=C2-B6-D03-NB-026,muc=NB,loigiai={Hai vectơ vuông góc khi tích vô hướng bằng $0$. Với $\overrightarrow{DA}$ và $\overrightarrow{GD}$ ta tính được tích vô hướng bằng $0$; các cặp còn lại có tích vô hướng khác $0$.}]{Cho hình hộp chữ nhật $ABCD.EFGH$ có $AB=3$, $AD=4$, $AE=4$. Cặp vectơ nào dưới đây vuông góc với nhau?}
  {$\overrightarrow{DA}$ và $\overrightarrow{GD}$.}
  {$\overrightarrow{BF}$ và $\overrightarrow{EB}$.}{$\overrightarrow{GE}$ và $\overrightarrow{EB}$.}{$\overrightarrow{GE}$ và $\overrightarrow{FC}$.}

\cauTN[id=C2-B6-D03-NB-027,muc=NB,loigiai={Hai vectơ vuông góc khi tích vô hướng bằng $0$. Với $\overrightarrow{H'E}$ và $\overrightarrow{FE}$ ta tính được tích vô hướng bằng $0$; các cặp còn lại có tích vô hướng khác $0$.}]{Cho hình lập phương $EFGH.E'F'G'H'$ có cạnh bằng $2$. Cặp vectơ nào dưới đây vuông góc với nhau?}
  {$\overrightarrow{H'E}$ và $\overrightarrow{FE}$.}
  {$\overrightarrow{EH}$ và $\overrightarrow{E'H}$.}{$\overrightarrow{F'E}$ và $\overrightarrow{FF'}$.}{$\overrightarrow{FH'}$ và $\overrightarrow{G'F}$.}

\cauTN[id=C2-B6-D03-NB-028,muc=NB,loigiai={Hai vectơ vuông góc khi tích vô hướng bằng $0$. Với $\overrightarrow{BF}$ và $\overrightarrow{DB}$ ta tính được tích vô hướng bằng $0$; các cặp còn lại có tích vô hướng khác $0$.}]{Cho hình hộp chữ nhật $ABCD.EFGH$ có $AB=2$, $AD=1$, $AE=4$. Cặp vectơ nào dưới đây vuông góc với nhau?}
  {$\overrightarrow{BF}$ và $\overrightarrow{DB}$.}
  {$\overrightarrow{GB}$ và $\overrightarrow{AF}$.}{$\overrightarrow{CA}$ và $\overrightarrow{CB}$.}{$\overrightarrow{DE}$ và $\overrightarrow{HF}$.}

\cauTN[id=C2-B6-D03-NB-029,muc=NB,loigiai={Hai vectơ vuông góc khi tích vô hướng bằng $0$. Với $\overrightarrow{MQ}$ và $\overrightarrow{NM'}$ ta tính được tích vô hướng bằng $0$; các cặp còn lại có tích vô hướng khác $0$.}]{Cho hình hộp chữ nhật $MNPQ.M'N'P'Q'$ có $MN=1$, $MQ=2$, $MM'=3$. Cặp vectơ nào dưới đây vuông góc với nhau?}
  {$\overrightarrow{MQ}$ và $\overrightarrow{NM'}$.}
  {$\overrightarrow{P'Q'}$ và $\overrightarrow{Q'N}$.}{$\overrightarrow{QN}$ và $\overrightarrow{QM}$.}{$\overrightarrow{M'N}$ và $\overrightarrow{M'N'}$.}

\cauTN[id=C2-B6-D03-NB-030,muc=NB,loigiai={Hai vectơ vuông góc khi tích vô hướng bằng $0$. Với $\overrightarrow{DG}$ và $\overrightarrow{AD}$ ta tính được tích vô hướng bằng $0$; các cặp còn lại có tích vô hướng khác $0$.}]{Cho hình hộp chữ nhật $ABCD.EFGH$ có $AB=2$, $AD=2$, $AE=1$. Cặp vectơ nào dưới đây vuông góc với nhau?}
  {$\overrightarrow{DG}$ và $\overrightarrow{AD}$.}
  {$\overrightarrow{EF}$ và $\overrightarrow{HC}$.}{$\overrightarrow{HF}$ và $\overrightarrow{FG}$.}{$\overrightarrow{DA}$ và $\overrightarrow{ED}$.}

\cauTN[id=C2-B6-D03-NB-031,muc=NB,loigiai={Hai vectơ vuông góc khi tích vô hướng bằng $0$. Với $\overrightarrow{Q'M'}$ và $\overrightarrow{M'N}$ ta tính được tích vô hướng bằng $0$; các cặp còn lại có tích vô hướng khác $0$.}]{Cho hình hộp chữ nhật $MNPQ.M'N'P'Q'$ có $MN=3$, $MQ=4$, $MM'=1$. Cặp vectơ nào dưới đây vuông góc với nhau?}
  {$\overrightarrow{Q'M'}$ và $\overrightarrow{M'N}$.}
  {$\overrightarrow{N'N}$ và $\overrightarrow{PM'}$.}{$\overrightarrow{PM'}$ và $\overrightarrow{MP'}$.}{$\overrightarrow{PM'}$ và $\overrightarrow{Q'M'}$.}

\cauTN[id=C2-B6-D03-NB-032,muc=NB,loigiai={Hai vectơ vuông góc khi tích vô hướng bằng $0$. Với $\overrightarrow{SC}$ và $\overrightarrow{BA}$ ta tính được tích vô hướng bằng $0$; các cặp còn lại có tích vô hướng khác $0$.}]{Cho tứ diện đều $SABC$ có cạnh bằng $2$. Cặp vectơ nào dưới đây vuông góc với nhau?}
  {$\overrightarrow{SC}$ và $\overrightarrow{BA}$.}
  {$\overrightarrow{CB}$ và $\overrightarrow{BA}$.}{$\overrightarrow{AC}$ và $\overrightarrow{BA}$.}{$\overrightarrow{CB}$ và $\overrightarrow{CS}$.}

\cauTN[id=C2-B6-D03-NB-033,muc=NB,loigiai={Hai vectơ vuông góc khi tích vô hướng bằng $0$. Với $\overrightarrow{N'Q'}$ và $\overrightarrow{PP'}$ ta tính được tích vô hướng bằng $0$; các cặp còn lại có tích vô hướng khác $0$.}]{Cho hình lập phương $MNPQ.M'N'P'Q'$ có cạnh bằng $a$. Cặp vectơ nào dưới đây vuông góc với nhau?}
  {$\overrightarrow{N'Q'}$ và $\overrightarrow{PP'}$.}
  {$\overrightarrow{NQ}$ và $\overrightarrow{N'P}$.}{$\overrightarrow{Q'N}$ và $\overrightarrow{N'Q}$.}{$\overrightarrow{M'N}$ và $\overrightarrow{PQ}$.}

\cauTN[id=C2-B6-D03-NB-034,muc=NB,loigiai={Hai vectơ vuông góc khi tích vô hướng bằng $0$. Với $\overrightarrow{BD}$ và $\overrightarrow{AC}$ ta tính được tích vô hướng bằng $0$; các cặp còn lại có tích vô hướng khác $0$.}]{Cho tứ diện đều $ABCD$ có cạnh bằng $a$. Cặp vectơ nào dưới đây vuông góc với nhau?}
  {$\overrightarrow{BD}$ và $\overrightarrow{AC}$.}
  {$\overrightarrow{CD}$ và $\overrightarrow{AD}$.}{$\overrightarrow{BD}$ và $\overrightarrow{DA}$.}{$\overrightarrow{AD}$ và $\overrightarrow{BD}$.}

\cauTN[id=C2-B6-D03-NB-035,muc=NB,loigiai={Hai vectơ vuông góc khi tích vô hướng bằng $0$. Với $\overrightarrow{F'G'}$ và $\overrightarrow{H'G}$ ta tính được tích vô hướng bằng $0$; các cặp còn lại có tích vô hướng khác $0$.}]{Cho hình lập phương $EFGH.E'F'G'H'$ có cạnh bằng $a$. Cặp vectơ nào dưới đây vuông góc với nhau?}
  {$\overrightarrow{F'G'}$ và $\overrightarrow{H'G}$.}
  {$\overrightarrow{G'G}$ và $\overrightarrow{F'H}$.}{$\overrightarrow{F'H'}$ và $\overrightarrow{E'H'}$.}{$\overrightarrow{HG'}$ và $\overrightarrow{E'G'}$.}

\cauTN[id=C2-B6-D03-NB-036,muc=NB,loigiai={Hai vectơ vuông góc khi tích vô hướng bằng $0$. Với $\overrightarrow{CB}$ và $\overrightarrow{SA}$ ta tính được tích vô hướng bằng $0$; các cặp còn lại có tích vô hướng khác $0$.}]{Cho tứ diện đều $SABC$ có cạnh bằng $3$. Cặp vectơ nào dưới đây vuông góc với nhau?}
  {$\overrightarrow{CB}$ và $\overrightarrow{SA}$.}
  {$\overrightarrow{BC}$ và $\overrightarrow{SB}$.}{$\overrightarrow{CS}$ và $\overrightarrow{AC}$.}{$\overrightarrow{AS}$ và $\overrightarrow{SC}$.}

\cauDS[id=C2-B6-D03-NB-037,muc=NB,loigiai={Tính các tích vô hướng qua ba cạnh xuất phát từ một đỉnh. a) Sai. b) Sai. c) Sai. d) Đúng.}]{Cho tứ diện đều $SABC$ có cạnh bằng $2$. Xét tính đúng sai của các mệnh đề sau:}{%
  \y{$\overrightarrow{CB}\cdot \overrightarrow{AC}=-1$.}{s}
  \y{$\overrightarrow{AB}\cdot \overrightarrow{AC}=1$.}{s}
  \y{$\overrightarrow{BS}\cdot \overrightarrow{AC}=2$.}{s}
  \y{$\overrightarrow{BC}\cdot \overrightarrow{BS}=2$.}{d}}

\cauDS[id=C2-B6-D03-NB-038,muc=NB,loigiai={Tính các tích vô hướng qua ba cạnh xuất phát từ một đỉnh. a) Sai. b) Sai. c) Sai. d) Đúng.}]{Cho hình lập phương $ABCD.EFGH$ có cạnh bằng $2$. Xét tính đúng sai của các mệnh đề sau:}{%
  \y{Hai vectơ $\overrightarrow{CB}$ và $\overrightarrow{BH}$ vuông góc với nhau.}{s}
  \y{$\overrightarrow{ED}\cdot \overrightarrow{GH}=1$.}{s}
  \y{$\overrightarrow{DA}\cdot \overrightarrow{CH}=-1$.}{s}
  \y{$\overrightarrow{GE}\cdot \overrightarrow{FB}=0$.}{d}}

\cauDS[id=C2-B6-D03-NB-039,muc=NB,loigiai={Tính các tích vô hướng qua ba cạnh xuất phát từ một đỉnh. a) Đúng. b) Sai. c) Đúng. d) Sai.}]{Cho tứ diện đều $SABC$ có cạnh bằng $1$. Xét tính đúng sai của các mệnh đề sau:}{%
  \y{$\overrightarrow{BA}\cdot \overrightarrow{SC}=0$.}{d}
  \y{$\overrightarrow{AC}\cdot \overrightarrow{AS}=\dfrac{3}{2}$.}{s}
  \y{$\overrightarrow{CA}\cdot \overrightarrow{AS}=-\dfrac{1}{2}$.}{d}
  \y{Hai vectơ $\overrightarrow{AC}$ và $\overrightarrow{SC}$ vuông góc với nhau.}{s}}

\cauDS[id=C2-B6-D03-NB-040,muc=NB,loigiai={Tính các tích vô hướng qua ba cạnh xuất phát từ một đỉnh. a) Sai. b) Đúng. c) Sai. d) Đúng.}]{Cho hình hộp chữ nhật $MNPQ.M'N'P'Q'$ có $MN=2$, $MQ=2$, $MM'=3$. Xét tính đúng sai của các mệnh đề sau:}{%
  \y{$\overrightarrow{M'Q}\cdot \overrightarrow{M'Q'}=-4$.}{s}
  \y{$\overrightarrow{P'P}\cdot \overrightarrow{P'Q'}=0$.}{d}
  \y{$\overrightarrow{QP}\cdot \overrightarrow{N'Q'}=-3$.}{s}
  \y{$\overrightarrow{QN}\cdot \overrightarrow{P'N}=4$.}{d}}

\cauDS[id=C2-B6-D03-NB-041,muc=NB,loigiai={Tính các tích vô hướng qua ba cạnh xuất phát từ một đỉnh. a) Đúng. b) Sai. c) Đúng. d) Sai.}]{Cho hình lập phương $EFGH.E'F'G'H'$ có cạnh bằng $2$. Xét tính đúng sai của các mệnh đề sau:}{%
  \y{$\overrightarrow{GF'}\cdot \overrightarrow{HG}=0$.}{d}
  \y{Hai vectơ $\overrightarrow{GH}$ và $\overrightarrow{H'F'}$ vuông góc với nhau.}{s}
  \y{Hai vectơ $\overrightarrow{H'F'}$ và $\overrightarrow{EG'}$ vuông góc với nhau.}{d}
  \y{Hai vectơ $\overrightarrow{FF'}$ và $\overrightarrow{F'G}$ vuông góc với nhau.}{s}}

\cauDS[id=C2-B6-D03-NB-042,muc=NB,loigiai={Tính các tích vô hướng qua ba cạnh xuất phát từ một đỉnh. a) Sai. b) Đúng. c) Sai. d) Đúng.}]{Cho tứ diện đều $ABCD$ có cạnh bằng $3$. Xét tính đúng sai của các mệnh đề sau:}{%
  \y{$\overrightarrow{CA}\cdot \overrightarrow{AC}=-8$.}{s}
  \y{$\overrightarrow{AB}\cdot \overrightarrow{CB}=\dfrac{9}{2}$.}{d}
  \y{$\overrightarrow{CD}\cdot \overrightarrow{AC}=-\dfrac{11}{2}$.}{s}
  \y{$\overrightarrow{CA}\cdot \overrightarrow{DC}=-\dfrac{9}{2}$.}{d}}

\cauDS[id=C2-B6-D03-NB-043,muc=NB,loigiai={Tính các tích vô hướng qua ba cạnh xuất phát từ một đỉnh. a) Sai. b) Sai. c) Đúng. d) Đúng.}]{Cho tứ diện đều $MNPQ$ có cạnh bằng $1$. Xét tính đúng sai của các mệnh đề sau:}{%
  \y{Hai vectơ $\overrightarrow{QN}$ và $\overrightarrow{QM}$ vuông góc với nhau.}{s}
  \y{$\overrightarrow{PQ}\cdot \overrightarrow{MQ}=1$.}{s}
  \y{$\overrightarrow{QM}\cdot \overrightarrow{NP}=0$.}{d}
  \y{$\overrightarrow{MP}\cdot \overrightarrow{QP}=\dfrac{1}{2}$.}{d}}

\cauDS[id=C2-B6-D03-NB-044,muc=NB,loigiai={Tính các tích vô hướng qua ba cạnh xuất phát từ một đỉnh. a) Đúng. b) Sai. c) Sai. d) Đúng.}]{Cho hình lập phương $MNPQ.M'N'P'Q'$ có cạnh bằng $2$. Xét tính đúng sai của các mệnh đề sau:}{%
  \y{$\overrightarrow{NM}\cdot \overrightarrow{Q'N}=-4$.}{d}
  \y{Hai vectơ $\overrightarrow{N'Q'}$ và $\overrightarrow{Q'P'}$ vuông góc với nhau.}{s}
  \y{$\overrightarrow{P'N}\cdot \overrightarrow{QQ'}=-3$.}{s}
  \y{$\overrightarrow{QM}\cdot \overrightarrow{NN'}=0$.}{d}}

\cauDS[id=C2-B6-D03-NB-045,muc=NB,loigiai={Tính các tích vô hướng qua ba cạnh xuất phát từ một đỉnh. a) Đúng. b) Sai. c) Đúng. d) Sai.}]{Cho tứ diện đều $MNPQ$ có cạnh bằng $3$. Xét tính đúng sai của các mệnh đề sau:}{%
  \y{$\overrightarrow{QP}\cdot \overrightarrow{MN}=0$.}{d}
  \y{$\overrightarrow{NP}\cdot \overrightarrow{MQ}=-1$.}{s}
  \y{$\overrightarrow{PN}\cdot \overrightarrow{NM}=-\dfrac{9}{2}$.}{d}
  \y{Hai vectơ $\overrightarrow{NM}$ và $\overrightarrow{MP}$ vuông góc với nhau.}{s}}

\cauDS[id=C2-B6-D03-NB-046,muc=NB,loigiai={Tính các tích vô hướng qua ba cạnh xuất phát từ một đỉnh. a) Sai. b) Sai. c) Đúng. d) Sai.}]{Cho hình hộp chữ nhật $MNPQ.M'N'P'Q'$ có $MN=2$, $MQ=3$, $MM'=4$. Xét tính đúng sai của các mệnh đề sau:}{%
  \y{$\overrightarrow{P'P}\cdot \overrightarrow{M'N}=17$.}{s}
  \y{Hai vectơ $\overrightarrow{PN}$ và $\overrightarrow{M'P}$ vuông góc với nhau.}{s}
  \y{$\overrightarrow{QP}\cdot \overrightarrow{PM}=-4$.}{d}
  \y{$\overrightarrow{NN'}\cdot \overrightarrow{QN}=2$.}{s}}

\cauDS[id=C2-B6-D03-NB-047,muc=NB,loigiai={Tính các tích vô hướng qua ba cạnh xuất phát từ một đỉnh. a) Đúng. b) Sai. c) Đúng. d) Sai.}]{Cho tứ diện đều $ABCD$ có cạnh bằng $a$. Xét tính đúng sai của các mệnh đề sau:}{%
  \y{$\overrightarrow{DC}\cdot \overrightarrow{BA}=0$.}{d}
  \y{Hai vectơ $\overrightarrow{DC}$ và $\overrightarrow{CA}$ vuông góc với nhau.}{s}
  \y{$\overrightarrow{BC}\cdot \overrightarrow{AD}=0$.}{d}
  \y{Hai vectơ $\overrightarrow{CD}$ và $\overrightarrow{DC}$ vuông góc với nhau.}{s}}

\cauDS[id=C2-B6-D03-NB-048,muc=NB,loigiai={Tính các tích vô hướng qua ba cạnh xuất phát từ một đỉnh. a) Đúng. b) Đúng. c) Đúng. d) Sai.}]{Cho hình lập phương $MNPQ.M'N'P'Q'$ có cạnh bằng $2$. Xét tính đúng sai của các mệnh đề sau:}{%
  \y{$\overrightarrow{NP'}\cdot \overrightarrow{Q'M'}=-4$.}{d}
  \y{$\overrightarrow{QN'}\cdot \overrightarrow{P'N}=0$.}{d}
  \y{$\overrightarrow{NQ'}\cdot \overrightarrow{NQ}=8$.}{d}
  \y{$\overrightarrow{MP}\cdot \overrightarrow{P'Q}=4$.}{s}}

\cauTLN[id=C2-B6-D03-NB-049,muc=NB,loigiai={Phân tích theo ba cạnh xuất phát từ một đỉnh ta được tích vô hướng bằng $-16$.}]{Cho hình hộp chữ nhật $EFGH.E'F'G'H'$ có $EF=4$, $EH=4$, $EE'=4$. Tính $\overrightarrow{E'G}\cdot \overrightarrow{FF'}$.}{-16}

\cauTLN[id=C2-B6-D03-NB-050,muc=NB,loigiai={Ta có $-\overrightarrow{DC}-\overrightarrow{AC}=-2\vec{b}+\vec{c}$ nên $|-\overrightarrow{DC}-\overrightarrow{AC}|^2=27$.}]{Cho tứ diện đều $ABCD$ có cạnh bằng $3$. Tính $|-\overrightarrow{DC}-\overrightarrow{AC}|^2$.}{27}

\cauTLN[id=C2-B6-D03-NB-051,muc=NB,loigiai={Phân tích theo ba cạnh xuất phát từ một đỉnh ta được tích vô hướng bằng $0$.}]{Cho hình lập phương $MNPQ.M'N'P'Q'$ có cạnh bằng $1$. Tính $\overrightarrow{PQ}\cdot \overrightarrow{PN}$.}{0}

\cauTLN[id=C2-B6-D03-NB-052,muc=NB,loigiai={Ta có $-\overrightarrow{QN'}=-\vec{a}+\vec{b}-\vec{c}$ nên $|-\overrightarrow{QN'}|^2=33$.}]{Cho hình hộp chữ nhật $MNPQ.M'N'P'Q'$ có $MN=4$, $MQ=1$, $MM'=4$. Tính $|-\overrightarrow{QN'}|^2$.}{33}

\cauTLN[id=C2-B6-D03-NB-053,muc=NB,loigiai={Ta có $\overrightarrow{B'B}+\overrightarrow{BD}=-\vec{a}+\vec{b}-\vec{c}$ nên $|\overrightarrow{B'B}+\overrightarrow{BD}|^2=12$.}]{Cho hình lập phương $ABCD.A'B'C'D'$ có cạnh bằng $2$. Tính $|\overrightarrow{B'B}+\overrightarrow{BD}|^2$.}{12}

\cauTLN[id=C2-B6-D03-NB-054,muc=NB,loigiai={Ta có $2\overrightarrow{GG'}+\overrightarrow{EH}=\vec{b}+2\vec{c}$ nên $|2\overrightarrow{GG'}+\overrightarrow{EH}|^2=13$.}]{Cho hình hộp chữ nhật $EFGH.E'F'G'H'$ có $EF=2$, $EH=3$, $EE'=1$. Tính $|2\overrightarrow{GG'}+\overrightarrow{EH}|^2$.}{13}

\cauTLN[id=C2-B6-D03-NB-055,muc=NB,loigiai={Ta có $2\overrightarrow{E'F}+2\overrightarrow{HG}=4\vec{a}-2\vec{c}$ nên $|2\overrightarrow{E'F}+2\overrightarrow{HG}|^2=20$.}]{Cho hình lập phương $EFGH.E'F'G'H'$ có cạnh bằng $1$. Tính $|2\overrightarrow{E'F}+2\overrightarrow{HG}|^2$.}{20}

\cauTLN[id=C2-B6-D03-NB-056,muc=NB,loigiai={Ta có $\overrightarrow{AC}+\overrightarrow{DB}=\vec{a}+\vec{b}-\vec{c}$ nên $|\overrightarrow{AC}+\overrightarrow{DB}|^2=8$.}]{Cho tứ diện đều $ABCD$ có cạnh bằng $2$. Tính $|\overrightarrow{AC}+\overrightarrow{DB}|^2$.}{8}

\cauTLN[id=C2-B6-D03-NB-057,muc=NB,loigiai={Phân tích theo ba cạnh xuất phát từ một đỉnh ta được tích vô hướng bằng $0$.}]{Cho tứ diện đều $SABC$ có cạnh bằng $1$. Tính $\overrightarrow{SC}\cdot \overrightarrow{BA}$.}{0}

\cauTLN[id=C2-B6-D03-NB-058,muc=NB,loigiai={Phân tích theo ba cạnh xuất phát từ một đỉnh ta được tích vô hướng bằng $0$.}]{Cho hình lập phương $ABCD.A'B'C'D'$ có cạnh bằng $3$. Tính $\overrightarrow{BB'}\cdot \overrightarrow{DC}$.}{0}

\cauTLN[id=C2-B6-D03-NB-059,muc=NB,loigiai={Phân tích theo ba cạnh xuất phát từ một đỉnh ta được tích vô hướng bằng $0$.}]{Cho hình lập phương $ABCD.A'B'C'D'$ có cạnh bằng $3$. Tính $\overrightarrow{AD}\cdot \overrightarrow{CD}$.}{0}

\cauTLN[id=C2-B6-D03-NB-060,muc=NB,loigiai={Phân tích theo ba cạnh xuất phát từ một đỉnh ta được tích vô hướng bằng $0$.}]{Cho hình lập phương $MNPQ.M'N'P'Q'$ có cạnh bằng $2$. Tính $\overrightarrow{M'M}\cdot \overrightarrow{NM}$.}{0}

\cauTN[id=C2-B6-D03-TH-001,muc=TH,loigiai={Biểu diễn hai vectơ theo ba cạnh xuất phát từ một đỉnh: $\overrightarrow{HG}=\vec{a}$, $\overrightarrow{AE}=\vec{c}$, rồi dùng $\vec{x}\cdot\vec{y}$ theo các tích vô hướng cơ bản ta được $0$.}]{Cho hình hộp chữ nhật $ABCD.EFGH$ có $AB=4$, $AD=1$, $AE=4$. Tích vô hướng $\overrightarrow{HG}\cdot \overrightarrow{AE}$ bằng}
  {$0$.}
  {$1$.}{$-1$.}{$2$.}

\cauTN[id=C2-B6-D03-TH-002,muc=TH,loigiai={Biểu diễn hai vectơ theo ba cạnh xuất phát từ một đỉnh: $\overrightarrow{PM}=-\vec{a}-\vec{b}$, $\overrightarrow{M'Q'}=\vec{b}$, rồi dùng $\vec{x}\cdot\vec{y}$ theo các tích vô hướng cơ bản ta được $-16$.}]{Cho hình hộp chữ nhật $MNPQ.M'N'P'Q'$ có $MN=1$, $MQ=4$, $MM'=2$. Tích vô hướng $\overrightarrow{PM}\cdot \overrightarrow{M'Q'}$ bằng}
  {$-16$.}
  {$-15$.}{$-17$.}{$16$.}

\cauTN[id=C2-B6-D03-TH-003,muc=TH,loigiai={Biểu diễn hai vectơ theo ba cạnh xuất phát từ một đỉnh: $\overrightarrow{H'G'}=\vec{a}$, $\overrightarrow{H'F}=\vec{a}-\vec{b}-\vec{c}$, rồi dùng $\vec{x}\cdot\vec{y}$ theo các tích vô hướng cơ bản ta được $16$.}]{Cho hình hộp chữ nhật $EFGH.E'F'G'H'$ có $EF=4$, $EH=2$, $EE'=3$. Tích vô hướng $\overrightarrow{H'G'}\cdot \overrightarrow{H'F}$ bằng}
  {$16$.}
  {$17$.}{$15$.}{$-16$.}

\cauTN[id=C2-B6-D03-TH-004,muc=TH,loigiai={Biểu diễn hai vectơ theo ba cạnh xuất phát từ một đỉnh: $\overrightarrow{QM'}=-\vec{b}+\vec{c}$, $\overrightarrow{N'P}=\vec{b}-\vec{c}$, rồi dùng $\vec{x}\cdot\vec{y}$ theo các tích vô hướng cơ bản ta được $-8$.}]{Cho hình lập phương $MNPQ.M'N'P'Q'$ có cạnh bằng $2$. Tích vô hướng $\overrightarrow{QM'}\cdot \overrightarrow{N'P}$ bằng}
  {$-8$.}
  {$-7$.}{$-9$.}{$8$.}

\cauTN[id=C2-B6-D03-TH-005,muc=TH,loigiai={Biểu diễn hai vectơ theo ba cạnh xuất phát từ một đỉnh: $\overrightarrow{B'C'}=\vec{b}$, $\overrightarrow{A'D'}=\vec{b}$, rồi dùng $\vec{x}\cdot\vec{y}$ theo các tích vô hướng cơ bản ta được $1$.}]{Cho hình hộp chữ nhật $ABCD.A'B'C'D'$ có $AB=1$, $AD=1$, $AA'=4$. Tích vô hướng $\overrightarrow{B'C'}\cdot \overrightarrow{A'D'}$ bằng}
  {$1$.}
  {$2$.}{$0$.}{$-1$.}

\cauTN[id=C2-B6-D03-TH-006,muc=TH,loigiai={Biểu diễn hai vectơ theo ba cạnh xuất phát từ một đỉnh: $\overrightarrow{DC'}=\vec{a}+\vec{c}$, $\overrightarrow{BB'}=\vec{c}$, rồi dùng $\vec{x}\cdot\vec{y}$ theo các tích vô hướng cơ bản ta được $4$.}]{Cho hình hộp chữ nhật $ABCD.A'B'C'D'$ có $AB=3$, $AD=3$, $AA'=2$. Tích vô hướng $\overrightarrow{DC'}\cdot \overrightarrow{BB'}$ bằng}
  {$4$.}
  {$5$.}{$3$.}{$-4$.}

\cauTN[id=C2-B6-D03-TH-007,muc=TH,loigiai={Biểu diễn hai vectơ theo ba cạnh xuất phát từ một đỉnh: $\overrightarrow{N'N}=-\vec{c}$, $\overrightarrow{P'P}=-\vec{c}$, rồi dùng $\vec{x}\cdot\vec{y}$ theo các tích vô hướng cơ bản ta được $1$.}]{Cho hình hộp chữ nhật $MNPQ.M'N'P'Q'$ có $MN=3$, $MQ=3$, $MM'=1$. Tích vô hướng $\overrightarrow{N'N}\cdot \overrightarrow{P'P}$ bằng}
  {$1$.}
  {$2$.}{$0$.}{$-1$.}

\cauTN[id=C2-B6-D03-TH-008,muc=TH,loigiai={Biểu diễn hai vectơ theo ba cạnh xuất phát từ một đỉnh: $\overrightarrow{BC}=-\vec{a}+\vec{b}$, $\overrightarrow{DB}=\vec{a}-\vec{c}$, rồi dùng $\vec{x}\cdot\vec{y}$ theo các tích vô hướng cơ bản ta được $-\dfrac{1}{2}$.}]{Cho tứ diện đều $ABCD$ có cạnh bằng $1$. Tích vô hướng $\overrightarrow{BC}\cdot \overrightarrow{DB}$ bằng}
  {$-\dfrac{1}{2}$.}
  {$\dfrac{1}{2}$.}{$-\dfrac{3}{2}$.}{$-1$.}

\cauTN[id=C2-B6-D03-TH-009,muc=TH,loigiai={Biểu diễn hai vectơ theo ba cạnh xuất phát từ một đỉnh: $\overrightarrow{QN}=\vec{a}-\vec{c}$, $\overrightarrow{PN}=\vec{a}-\vec{b}$, rồi dùng $\vec{x}\cdot\vec{y}$ theo các tích vô hướng cơ bản ta được $\dfrac{9}{2}$.}]{Cho tứ diện đều $MNPQ$ có cạnh bằng $3$. Tích vô hướng $\overrightarrow{QN}\cdot \overrightarrow{PN}$ bằng}
  {$\dfrac{9}{2}$.}
  {$\dfrac{11}{2}$.}{$\dfrac{7}{2}$.}{$-\dfrac{9}{2}$.}

\cauTN[id=C2-B6-D03-TH-010,muc=TH,loigiai={Biểu diễn hai vectơ theo ba cạnh xuất phát từ một đỉnh: $\overrightarrow{DB}=\vec{a}-\vec{c}$, $\overrightarrow{BC}=-\vec{a}+\vec{b}$, rồi dùng $\vec{x}\cdot\vec{y}$ theo các tích vô hướng cơ bản ta được $-\dfrac{1}{2}a^2$.}]{Cho tứ diện đều $ABCD$ có cạnh bằng $a$. Tích vô hướng $\overrightarrow{DB}\cdot \overrightarrow{BC}$ bằng}
  {$-\dfrac{1}{2}a^2$.}
  {$\dfrac{1}{2}a^2$.}{$-\dfrac{3}{2}a^2$.}{$-a^2$.}

\cauTN[id=C2-B6-D03-TH-011,muc=TH,loigiai={Biểu diễn hai vectơ theo ba cạnh xuất phát từ một đỉnh: $\overrightarrow{DA}=-\vec{b}$, $\overrightarrow{DB'}=\vec{a}-\vec{b}+\vec{c}$, rồi dùng $\vec{x}\cdot\vec{y}$ theo các tích vô hướng cơ bản ta được $1$.}]{Cho hình lập phương $ABCD.A'B'C'D'$ có cạnh bằng $1$. Tích vô hướng $\overrightarrow{DA}\cdot \overrightarrow{DB'}$ bằng}
  {$1$.}
  {$2$.}{$0$.}{$-1$.}

\cauTN[id=C2-B6-D03-TH-012,muc=TH,loigiai={Biểu diễn hai vectơ theo ba cạnh xuất phát từ một đỉnh: $\overrightarrow{HF}=\vec{a}-\vec{b}$, $\overrightarrow{EA}=-\vec{c}$, rồi dùng $\vec{x}\cdot\vec{y}$ theo các tích vô hướng cơ bản ta được $0$.}]{Cho hình hộp chữ nhật $ABCD.EFGH$ có $AB=3$, $AD=4$, $AE=2$. Tích vô hướng $\overrightarrow{HF}\cdot \overrightarrow{EA}$ bằng}
  {$0$.}
  {$1$.}{$-1$.}{$2$.}

\cauTN[id=C2-B6-D03-TH-013,muc=TH,loigiai={Ta có $\overrightarrow{MP}\cdot \overrightarrow{NM}=-2$, $|\overrightarrow{MP}|^2=4$, $|\overrightarrow{NM}|^2=4$ nên $\cos=\dfrac{-2}{\sqrt{4\cdot 4}}$, suy ra góc bằng $120^\circ$.}]{Cho tứ diện đều $MNPQ$ có cạnh bằng $2$. Góc giữa hai vectơ $\overrightarrow{MP}$ và $\overrightarrow{NM}$ bằng}
  {$120^\circ$.}
  {$90^\circ$.}{$30^\circ$.}{$150^\circ$.}

\cauTN[id=C2-B6-D03-TH-014,muc=TH,loigiai={Ta có $\overrightarrow{B'C}\cdot \overrightarrow{BC'}=0$, $|\overrightarrow{B'C}|^2=2a^2$, $|\overrightarrow{BC'}|^2=2a^2$ nên $\cos=\dfrac{0}{\sqrt{2a^2\cdot 2a^2}}$, suy ra góc bằng $90^\circ$.}]{Cho hình lập phương $ABCD.A'B'C'D'$ có cạnh bằng $a$. Góc giữa hai vectơ $\overrightarrow{B'C}$ và $\overrightarrow{BC'}$ bằng}
  {$90^\circ$.}
  {$150^\circ$.}{$30^\circ$.}{$120^\circ$.}

\cauTN[id=C2-B6-D03-TH-015,muc=TH,loigiai={Ta có $\overrightarrow{AC}\cdot \overrightarrow{AS}=2$, $|\overrightarrow{AC}|^2=4$, $|\overrightarrow{AS}|^2=4$ nên $\cos=\dfrac{2}{\sqrt{4\cdot 4}}$, suy ra góc bằng $60^\circ$.}]{Cho tứ diện đều $SABC$ có cạnh bằng $2$. Góc giữa hai vectơ $\overrightarrow{AC}$ và $\overrightarrow{AS}$ bằng}
  {$60^\circ$.}
  {$90^\circ$.}{$45^\circ$.}{$150^\circ$.}

\cauTN[id=C2-B6-D03-TH-016,muc=TH,loigiai={Ta có $\overrightarrow{NM}\cdot \overrightarrow{QM}=\dfrac{1}{2}$, $|\overrightarrow{NM}|^2=1$, $|\overrightarrow{QM}|^2=1$ nên $\cos=\dfrac{\dfrac{1}{2}}{\sqrt{1\cdot 1}}$, suy ra góc bằng $60^\circ$.}]{Cho tứ diện đều $MNPQ$ có cạnh bằng $1$. Góc giữa hai vectơ $\overrightarrow{NM}$ và $\overrightarrow{QM}$ bằng}
  {$60^\circ$.}
  {$45^\circ$.}{$150^\circ$.}{$135^\circ$.}

\cauTN[id=C2-B6-D03-TH-017,muc=TH,loigiai={Ta có $\overrightarrow{CA}\cdot \overrightarrow{DD'}=0$, $|\overrightarrow{CA}|^2=2a^2$, $|\overrightarrow{DD'}|^2=a^2$ nên $\cos=\dfrac{0}{\sqrt{2a^2\cdot a^2}}$, suy ra góc bằng $90^\circ$.}]{Cho hình lập phương $ABCD.A'B'C'D'$ có cạnh bằng $a$. Góc giữa hai vectơ $\overrightarrow{CA}$ và $\overrightarrow{DD'}$ bằng}
  {$90^\circ$.}
  {$135^\circ$.}{$45^\circ$.}{$120^\circ$.}

\cauTN[id=C2-B6-D03-TH-018,muc=TH,loigiai={Ta có $\overrightarrow{AF}\cdot \overrightarrow{EB}=0$, $|\overrightarrow{AF}|^2=18$, $|\overrightarrow{EB}|^2=18$ nên $\cos=\dfrac{0}{\sqrt{18\cdot 18}}$, suy ra góc bằng $90^\circ$.}]{Cho hình lập phương $ABCD.EFGH$ có cạnh bằng $3$. Góc giữa hai vectơ $\overrightarrow{AF}$ và $\overrightarrow{EB}$ bằng}
  {$90^\circ$.}
  {$150^\circ$.}{$120^\circ$.}{$45^\circ$.}

\cauTN[id=C2-B6-D03-TH-019,muc=TH,loigiai={Ta có $\overrightarrow{DA}\cdot \overrightarrow{AC}=-\dfrac{1}{2}$, $|\overrightarrow{DA}|^2=1$, $|\overrightarrow{AC}|^2=1$ nên $\cos=\dfrac{-\dfrac{1}{2}}{\sqrt{1\cdot 1}}$, suy ra góc bằng $120^\circ$.}]{Cho tứ diện đều $ABCD$ có cạnh bằng $1$. Góc giữa hai vectơ $\overrightarrow{DA}$ và $\overrightarrow{AC}$ bằng}
  {$120^\circ$.}
  {$90^\circ$.}{$60^\circ$.}{$30^\circ$.}

\cauTN[id=C2-B6-D03-TH-020,muc=TH,loigiai={Ta có $\overrightarrow{FC}\cdot \overrightarrow{CA}=-4$, $|\overrightarrow{FC}|^2=8$, $|\overrightarrow{CA}|^2=8$ nên $\cos=\dfrac{-4}{\sqrt{8\cdot 8}}$, suy ra góc bằng $120^\circ$.}]{Cho hình lập phương $ABCD.EFGH$ có cạnh bằng $2$. Góc giữa hai vectơ $\overrightarrow{FC}$ và $\overrightarrow{CA}$ bằng}
  {$120^\circ$.}
  {$30^\circ$.}{$135^\circ$.}{$90^\circ$.}

\cauTN[id=C2-B6-D03-TH-021,muc=TH,loigiai={Ta có $\overrightarrow{DB}\cdot \overrightarrow{AC}=0$, $|\overrightarrow{DB}|^2=9$, $|\overrightarrow{AC}|^2=9$ nên $\cos=\dfrac{0}{\sqrt{9\cdot 9}}$, suy ra góc bằng $90^\circ$.}]{Cho tứ diện đều $ABCD$ có cạnh bằng $3$. Góc giữa hai vectơ $\overrightarrow{DB}$ và $\overrightarrow{AC}$ bằng}
  {$90^\circ$.}
  {$60^\circ$.}{$120^\circ$.}{$30^\circ$.}

\cauTN[id=C2-B6-D03-TH-022,muc=TH,loigiai={Ta có $\overrightarrow{AB}\cdot \overrightarrow{BC}=-2$, $|\overrightarrow{AB}|^2=4$, $|\overrightarrow{BC}|^2=4$ nên $\cos=\dfrac{-2}{\sqrt{4\cdot 4}}$, suy ra góc bằng $120^\circ$.}]{Cho tứ diện đều $ABCD$ có cạnh bằng $2$. Góc giữa hai vectơ $\overrightarrow{AB}$ và $\overrightarrow{BC}$ bằng}
  {$120^\circ$.}
  {$135^\circ$.}{$60^\circ$.}{$30^\circ$.}

\cauTN[id=C2-B6-D03-TH-023,muc=TH,loigiai={Ta có $\overrightarrow{MP}\cdot \overrightarrow{PQ}=-4$, $|\overrightarrow{MP}|^2=8$, $|\overrightarrow{PQ}|^2=4$ nên $\cos=\dfrac{-4}{\sqrt{8\cdot 4}}$, suy ra góc bằng $135^\circ$.}]{Cho hình lập phương $MNPQ.M'N'P'Q'$ có cạnh bằng $2$. Góc giữa hai vectơ $\overrightarrow{MP}$ và $\overrightarrow{PQ}$ bằng}
  {$135^\circ$.}
  {$90^\circ$.}{$30^\circ$.}{$45^\circ$.}

\cauTN[id=C2-B6-D03-TH-024,muc=TH,loigiai={Ta có $\overrightarrow{HD}\cdot \overrightarrow{BC}=0$, $|\overrightarrow{HD}|^2=1$, $|\overrightarrow{BC}|^2=1$ nên $\cos=\dfrac{0}{\sqrt{1\cdot 1}}$, suy ra góc bằng $90^\circ$.}]{Cho hình lập phương $ABCD.EFGH$ có cạnh bằng $1$. Góc giữa hai vectơ $\overrightarrow{HD}$ và $\overrightarrow{BC}$ bằng}
  {$90^\circ$.}
  {$135^\circ$.}{$60^\circ$.}{$150^\circ$.}

\cauTN[id=C2-B6-D03-TH-025,muc=TH,loigiai={Ta có $\vec{a}\cdot\vec{b}=1\cdot3\cdot\cos 120^\circ=-\dfrac{3}{2}$. Khai triển: $(-\vec{a}+2\vec{b})\cdot(\vec{a}+2\vec{b})=-1\vec{a}^2+0\vec{a}\cdot\vec{b}+4\vec{b}^2=35$.}]{Cho hai vectơ $\vec{a}$, $\vec{b}$ có $|\vec{a}|=1$, $|\vec{b}|=3$ và $(\vec{a},\vec{b})=120^\circ$. Tích vô hướng $(-\vec{a}+2\vec{b})\cdot(\vec{a}+2\vec{b})$ bằng}
  {$35$.}
  {$33$.}{$71$.}{$36$.}

\cauTN[id=C2-B6-D03-TH-026,muc=TH,loigiai={Ta có $\vec{a}\cdot\vec{b}=1\cdot3\cdot\cos 90^\circ=0$. Khai triển: $(\vec{a}-\vec{b})\cdot(\vec{a}-\vec{b})=1\vec{a}^2-2\vec{a}\cdot\vec{b}+1\vec{b}^2=10$.}]{Cho hai vectơ $\vec{a}$, $\vec{b}$ có $|\vec{a}|=1$, $|\vec{b}|=3$ và $(\vec{a},\vec{b})=90^\circ$. Tích vô hướng $(\vec{a}-\vec{b})\cdot(\vec{a}-\vec{b})$ bằng}
  {$10$.}
  {$21$.}{$20$.}{$5$.}

\cauTN[id=C2-B6-D03-TH-027,muc=TH,loigiai={Ta có $\vec{a}\cdot\vec{b}=2\cdot2\cdot\cos 120^\circ=-2$. Khai triển: $(-\vec{a}-\vec{b})\cdot(-\vec{a}+2\vec{b})=1\vec{a}^2-1\vec{a}\cdot\vec{b}-2\vec{b}^2=-2$.}]{Cho hai vectơ $\vec{a}$, $\vec{b}$ có $|\vec{a}|=2$, $|\vec{b}|=2$ và $(\vec{a},\vec{b})=120^\circ$. Tích vô hướng $(-\vec{a}-\vec{b})\cdot(-\vec{a}+2\vec{b})$ bằng}
  {$-2$.}
  {$-4$.}{$1$.}{$-1$.}

\cauTN[id=C2-B6-D03-TH-028,muc=TH,loigiai={Ta có $\vec{a}\cdot\vec{b}=4\cdot1\cdot\cos 90^\circ=0$. Khai triển: $(-\vec{a}+3\vec{b})\cdot(2\vec{a}+3\vec{b})=-2\vec{a}^2+3\vec{a}\cdot\vec{b}+9\vec{b}^2=-23$.}]{Cho hai vectơ $\vec{a}$, $\vec{b}$ có $|\vec{a}|=4$, $|\vec{b}|=1$ và $(\vec{a},\vec{b})=90^\circ$. Tích vô hướng $(-\vec{a}+3\vec{b})\cdot(2\vec{a}+3\vec{b})$ bằng}
  {$-23$.}
  {$-24$.}{$-\dfrac{23}{2}$.}{$-46$.}

\cauTN[id=C2-B6-D03-TH-029,muc=TH,loigiai={Ta có $\vec{a}\cdot\vec{b}=1\cdot5\cdot\cos 60^\circ=\dfrac{5}{2}$. Khai triển: $(2\vec{a}+2\vec{b})\cdot(3\vec{a}+\vec{b})=6\vec{a}^2+8\vec{a}\cdot\vec{b}+2\vec{b}^2=76$.}]{Cho hai vectơ $\vec{a}$, $\vec{b}$ có $|\vec{a}|=1$, $|\vec{b}|=5$ và $(\vec{a},\vec{b})=60^\circ$. Tích vô hướng $(2\vec{a}+2\vec{b})\cdot(3\vec{a}+\vec{b})$ bằng}
  {$76$.}
  {$56$.}{$38$.}{$78$.}

\cauTN[id=C2-B6-D03-TH-030,muc=TH,loigiai={Ta có $\vec{a}\cdot\vec{b}=1\cdot3\cdot\cos 90^\circ=0$. Khai triển: $(\vec{a}+3\vec{b})\cdot(2\vec{a}+2\vec{b})=2\vec{a}^2+8\vec{a}\cdot\vec{b}+6\vec{b}^2=56$.}]{Cho hai vectơ $\vec{a}$, $\vec{b}$ có $|\vec{a}|=1$, $|\vec{b}|=3$ và $(\vec{a},\vec{b})=90^\circ$. Tích vô hướng $(\vec{a}+3\vec{b})\cdot(2\vec{a}+2\vec{b})$ bằng}
  {$56$.}
  {$53$.}{$113$.}{$168$.}

\cauTN[id=C2-B6-D03-TH-031,muc=TH,loigiai={Ta có $\vec{a}\cdot\vec{b}=1\cdot5\cdot\cos 60^\circ=\dfrac{5}{2}$. Khai triển: $(-\vec{a}+3\vec{b})\cdot(-\vec{a}+3\vec{b})=1\vec{a}^2-6\vec{a}\cdot\vec{b}+9\vec{b}^2=211$.}]{Cho hai vectơ $\vec{a}$, $\vec{b}$ có $|\vec{a}|=1$, $|\vec{b}|=5$ và $(\vec{a},\vec{b})=60^\circ$. Tích vô hướng $(-\vec{a}+3\vec{b})\cdot(-\vec{a}+3\vec{b})$ bằng}
  {$211$.}
  {$423$.}{$226$.}{$633$.}

\cauTN[id=C2-B6-D03-TH-032,muc=TH,loigiai={Ta có $\vec{a}\cdot\vec{b}=4\cdot3\cdot\cos 60^\circ=6$. Khai triển: $(-\vec{a}+2\vec{b})\cdot(\vec{a}-\vec{b})=-1\vec{a}^2+3\vec{a}\cdot\vec{b}-2\vec{b}^2=-16$.}]{Cho hai vectơ $\vec{a}$, $\vec{b}$ có $|\vec{a}|=4$, $|\vec{b}|=3$ và $(\vec{a},\vec{b})=60^\circ$. Tích vô hướng $(-\vec{a}+2\vec{b})\cdot(\vec{a}-\vec{b})$ bằng}
  {$-16$.}
  {$-18$.}{$-48$.}{$-17$.}

\cauTN[id=C2-B6-D03-TH-033,muc=TH,loigiai={Ta có $\vec{a}\cdot\vec{b}=1\cdot1\cdot\cos 120^\circ=-\dfrac{1}{2}$. Khai triển: $(2\vec{a}+2\vec{b})\cdot(-\vec{a}-\vec{b})=-2\vec{a}^2-4\vec{a}\cdot\vec{b}-2\vec{b}^2=-2$.}]{Cho hai vectơ $\vec{a}$, $\vec{b}$ có $|\vec{a}|=1$, $|\vec{b}|=1$ và $(\vec{a},\vec{b})=120^\circ$. Tích vô hướng $(2\vec{a}+2\vec{b})\cdot(-\vec{a}-\vec{b})$ bằng}
  {$-2$.}
  {$-4$.}{$1$.}{$2$.}

\cauTN[id=C2-B6-D03-TH-034,muc=TH,loigiai={Ta có $\vec{a}\cdot\vec{b}=3\cdot2\cdot\cos 120^\circ=-3$. Khai triển: $(-\vec{a}+2\vec{b})\cdot(\vec{a}+\vec{b})=-1\vec{a}^2+1\vec{a}\cdot\vec{b}+2\vec{b}^2=-4$.}]{Cho hai vectơ $\vec{a}$, $\vec{b}$ có $|\vec{a}|=3$, $|\vec{b}|=2$ và $(\vec{a},\vec{b})=120^\circ$. Tích vô hướng $(-\vec{a}+2\vec{b})\cdot(\vec{a}+\vec{b})$ bằng}
  {$-4$.}
  {$-1$.}{$-12$.}{$-5$.}

\cauTN[id=C2-B6-D03-TH-035,muc=TH,loigiai={Ta có $\vec{a}\cdot\vec{b}=4\cdot5\cdot\cos 60^\circ=10$. Khai triển: $(-\vec{a}+3\vec{b})\cdot(2\vec{a}+\vec{b})=-2\vec{a}^2+5\vec{a}\cdot\vec{b}+3\vec{b}^2=93$.}]{Cho hai vectơ $\vec{a}$, $\vec{b}$ có $|\vec{a}|=4$, $|\vec{b}|=5$ và $(\vec{a},\vec{b})=60^\circ$. Tích vô hướng $(-\vec{a}+3\vec{b})\cdot(2\vec{a}+\vec{b})$ bằng}
  {$93$.}
  {$92$.}{$96$.}{$187$.}

\cauTN[id=C2-B6-D03-TH-036,muc=TH,loigiai={Ta có $\vec{a}\cdot\vec{b}=1\cdot5\cdot\cos 60^\circ=\dfrac{5}{2}$. Khai triển: $(-\vec{a}-\vec{b})\cdot(\vec{a}+3\vec{b})=-1\vec{a}^2-4\vec{a}\cdot\vec{b}-3\vec{b}^2=-86$.}]{Cho hai vectơ $\vec{a}$, $\vec{b}$ có $|\vec{a}|=1$, $|\vec{b}|=5$ và $(\vec{a},\vec{b})=60^\circ$. Tích vô hướng $(-\vec{a}-\vec{b})\cdot(\vec{a}+3\vec{b})$ bằng}
  {$-86$.}
  {$-172$.}{$-258$.}{$-43$.}

\cauTN[id=C2-B6-D03-TH-037,muc=TH,loigiai={Biểu diễn $2\overrightarrow{QP}+2\overrightarrow{PM'}=-2\vec{b}+2\vec{c}$ theo ba cạnh xuất phát từ một đỉnh rồi dùng $|\vec{x}|^2=\vec{x}\cdot\vec{x}$ ta được $|2\overrightarrow{QP}+2\overrightarrow{PM'}|^2=52$, nên $|2\overrightarrow{QP}+2\overrightarrow{PM'}|=2\sqrt{13}$.}]{Cho hình hộp chữ nhật $MNPQ.M'N'P'Q'$ có $MN=2$, $MQ=3$, $MM'=2$. Độ dài của vectơ $2\overrightarrow{QP}+2\overrightarrow{PM'}$ bằng}
  {$2\sqrt{13}$.}
  {$\sqrt{17}$.}{$2\sqrt{14}$.}{$2\sqrt{29}$.}

\cauTN[id=C2-B6-D03-TH-038,muc=TH,loigiai={Biểu diễn $-\overrightarrow{Q'N}-\overrightarrow{Q'M}=-\vec{a}+2\vec{b}+2\vec{c}$ theo ba cạnh xuất phát từ một đỉnh rồi dùng $|\vec{x}|^2=\vec{x}\cdot\vec{x}$ ta được $|-\overrightarrow{Q'N}-\overrightarrow{Q'M}|^2=101$, nên $|-\overrightarrow{Q'N}-\overrightarrow{Q'M}|=\sqrt{101}$.}]{Cho hình hộp chữ nhật $MNPQ.M'N'P'Q'$ có $MN=1$, $MQ=4$, $MM'=3$. Độ dài của vectơ $-\overrightarrow{Q'N}-\overrightarrow{Q'M}$ bằng}
  {$\sqrt{101}$.}
  {$\sqrt{229}$.}{$1$.}{$\sqrt{226}$.}

\cauTN[id=C2-B6-D03-TH-039,muc=TH,loigiai={Biểu diễn $2\overrightarrow{E'H'}+2\overrightarrow{G'E'}=-2\vec{a}$ theo ba cạnh xuất phát từ một đỉnh rồi dùng $|\vec{x}|^2=\vec{x}\cdot\vec{x}$ ta được $|2\overrightarrow{E'H'}+2\overrightarrow{G'E'}|^2=4$, nên $|2\overrightarrow{E'H'}+2\overrightarrow{G'E'}|=2$.}]{Cho hình lập phương $EFGH.E'F'G'H'$ có cạnh bằng $1$. Độ dài của vectơ $2\overrightarrow{E'H'}+2\overrightarrow{G'E'}$ bằng}
  {$2$.}
  {$\sqrt{2}$.}{$2\sqrt{5}$.}{$\sqrt{5}$.}

\cauTN[id=C2-B6-D03-TH-040,muc=TH,loigiai={Biểu diễn $\overrightarrow{Q'N}+\overrightarrow{P'N}=\vec{a}-2\vec{b}-2\vec{c}$ theo ba cạnh xuất phát từ một đỉnh rồi dùng $|\vec{x}|^2=\vec{x}\cdot\vec{x}$ ta được $|\overrightarrow{Q'N}+\overrightarrow{P'N}|^2=104$, nên $|\overrightarrow{Q'N}+\overrightarrow{P'N}|=2\sqrt{26}$.}]{Cho hình hộp chữ nhật $MNPQ.M'N'P'Q'$ có $MN=2$, $MQ=4$, $MM'=3$. Độ dài của vectơ $\overrightarrow{Q'N}+\overrightarrow{P'N}$ bằng}
  {$2\sqrt{26}$.}
  {$2$.}{$\sqrt{229}$.}{$\sqrt{241}$.}

\cauTN[id=C2-B6-D03-TH-041,muc=TH,loigiai={Biểu diễn $2\overrightarrow{A'D'}-\overrightarrow{AC}=-\vec{a}+\vec{b}$ theo ba cạnh xuất phát từ một đỉnh rồi dùng $|\vec{x}|^2=\vec{x}\cdot\vec{x}$ ta được $|2\overrightarrow{A'D'}-\overrightarrow{AC}|^2=10$, nên $|2\overrightarrow{A'D'}-\overrightarrow{AC}|=\sqrt{10}$.}]{Cho hình hộp chữ nhật $ABCD.A'B'C'D'$ có $AB=1$, $AD=3$, $AA'=3$. Độ dài của vectơ $2\overrightarrow{A'D'}-\overrightarrow{AC}$ bằng}
  {$\sqrt{10}$.}
  {$\sqrt{82}$.}{$1$.}{$2$.}

\cauTN[id=C2-B6-D03-TH-042,muc=TH,loigiai={Biểu diễn $-\overrightarrow{CB}+\overrightarrow{BG}=2\vec{b}+\vec{c}$ theo ba cạnh xuất phát từ một đỉnh rồi dùng $|\vec{x}|^2=\vec{x}\cdot\vec{x}$ ta được $|-\overrightarrow{CB}+\overrightarrow{BG}|^2=73$, nên $|-\overrightarrow{CB}+\overrightarrow{BG}|=\sqrt{73}$.}]{Cho hình hộp chữ nhật $ABCD.EFGH$ có $AB=1$, $AD=4$, $AE=3$. Độ dài của vectơ $-\overrightarrow{CB}+\overrightarrow{BG}$ bằng}
  {$\sqrt{73}$.}
  {$6\sqrt{5}$.}{$3$.}{$3\sqrt{17}$.}

\cauTN[id=C2-B6-D03-TH-043,muc=TH,loigiai={Biểu diễn $\overrightarrow{BC}+2\overrightarrow{AB}=\vec{a}+\vec{b}$ theo ba cạnh xuất phát từ một đỉnh rồi dùng $|\vec{x}|^2=\vec{x}\cdot\vec{x}$ ta được $|\overrightarrow{BC}+2\overrightarrow{AB}|^2=27$, nên $|\overrightarrow{BC}+2\overrightarrow{AB}|=3\sqrt{3}$.}]{Cho tứ diện đều $ABCD$ có cạnh bằng $3$. Độ dài của vectơ $\overrightarrow{BC}+2\overrightarrow{AB}$ bằng}
  {$3\sqrt{3}$.}
  {$6$.}{$3\sqrt{7}$.}{$3$.}

\cauTN[id=C2-B6-D03-TH-044,muc=TH,loigiai={Biểu diễn $2\overrightarrow{F'E'}+\overrightarrow{FH}=-3\vec{a}+\vec{b}$ theo ba cạnh xuất phát từ một đỉnh rồi dùng $|\vec{x}|^2=\vec{x}\cdot\vec{x}$ ta được $|2\overrightarrow{F'E'}+\overrightarrow{FH}|^2=45$, nên $|2\overrightarrow{F'E'}+\overrightarrow{FH}|=3\sqrt{5}$.}]{Cho hình hộp chữ nhật $EFGH.E'F'G'H'$ có $EF=2$, $EH=3$, $EE'=4$. Độ dài của vectơ $2\overrightarrow{F'E'}+\overrightarrow{FH}$ bằng}
  {$3\sqrt{5}$.}
  {$5$.}{$10$.}{$\sqrt{13}$.}

\cauTN[id=C2-B6-D03-TH-045,muc=TH,loigiai={Biểu diễn $2\overrightarrow{C'A}+\overrightarrow{AD}=-2\vec{a}-\vec{b}-2\vec{c}$ theo ba cạnh xuất phát từ một đỉnh rồi dùng $|\vec{x}|^2=\vec{x}\cdot\vec{x}$ ta được $|2\overrightarrow{C'A}+\overrightarrow{AD}|^2=33$, nên $|2\overrightarrow{C'A}+\overrightarrow{AD}|=\sqrt{33}$.}]{Cho hình hộp chữ nhật $ABCD.A'B'C'D'$ có $AB=2$, $AD=1$, $AA'=2$. Độ dài của vectơ $2\overrightarrow{C'A}+\overrightarrow{AD}$ bằng}
  {$\sqrt{33}$.}
  {$\sqrt{41}$.}{$4\sqrt{2}$.}{$2\sqrt{2}$.}

\cauTN[id=C2-B6-D03-TH-046,muc=TH,loigiai={Biểu diễn $-\overrightarrow{FA}+\overrightarrow{AE}=\vec{a}+2\vec{c}$ theo ba cạnh xuất phát từ một đỉnh rồi dùng $|\vec{x}|^2=\vec{x}\cdot\vec{x}$ ta được $|-\overrightarrow{FA}+\overrightarrow{AE}|^2=65$, nên $|-\overrightarrow{FA}+\overrightarrow{AE}|=\sqrt{65}$.}]{Cho hình hộp chữ nhật $ABCD.EFGH$ có $AB=1$, $AD=2$, $AE=4$. Độ dài của vectơ $-\overrightarrow{FA}+\overrightarrow{AE}$ bằng}
  {$\sqrt{65}$.}
  {$2\sqrt{37}$.}{$1$.}{$\sqrt{145}$.}

\cauTN[id=C2-B6-D03-TH-047,muc=TH,loigiai={Biểu diễn $2\overrightarrow{CA}-\overrightarrow{DA}=-2\vec{b}+\vec{c}$ theo ba cạnh xuất phát từ một đỉnh rồi dùng $|\vec{x}|^2=\vec{x}\cdot\vec{x}$ ta được $|2\overrightarrow{CA}-\overrightarrow{DA}|^2=12$, nên $|2\overrightarrow{CA}-\overrightarrow{DA}|=2\sqrt{3}$.}]{Cho tứ diện đều $ABCD$ có cạnh bằng $2$. Độ dài của vectơ $2\overrightarrow{CA}-\overrightarrow{DA}$ bằng}
  {$2\sqrt{3}$.}
  {$2\sqrt{7}$.}{$2$.}{$4$.}

\cauTN[id=C2-B6-D03-TH-048,muc=TH,loigiai={Biểu diễn $\overrightarrow{B'A}+2\overrightarrow{DA'}=-\vec{a}-2\vec{b}+\vec{c}$ theo ba cạnh xuất phát từ một đỉnh rồi dùng $|\vec{x}|^2=\vec{x}\cdot\vec{x}$ ta được $|\overrightarrow{B'A}+2\overrightarrow{DA'}|^2=17$, nên $|\overrightarrow{B'A}+2\overrightarrow{DA'}|=\sqrt{17}$.}]{Cho hình hộp chữ nhật $ABCD.A'B'C'D'$ có $AB=2$, $AD=1$, $AA'=3$. Độ dài của vectơ $\overrightarrow{B'A}+2\overrightarrow{DA'}$ bằng}
  {$\sqrt{17}$.}
  {$\sqrt{89}$.}{$2\sqrt{5}$.}{$\sqrt{5}$.}

\cauDS[id=C2-B6-D03-TH-049,muc=TH,loigiai={Tính các tích vô hướng qua ba cạnh xuất phát từ một đỉnh. a) Sai. b) Sai. c) Đúng. d) Đúng.}]{Cho hình lập phương $EFGH.E'F'G'H'$ có cạnh bằng $a$. Xét tính đúng sai của các mệnh đề sau:}{%
  \y{$|\overrightarrow{H'H}|^2=-a^2$.}{s}
  \y{$|\overrightarrow{HE}|^2=0$.}{s}
  \y{$\overrightarrow{H'E'}\cdot \overrightarrow{FE}=0$.}{d}
  \y{$\overrightarrow{HE}\cdot \overrightarrow{EG}=-a^2$.}{d}}

\cauDS[id=C2-B6-D03-TH-050,muc=TH,loigiai={Tính các tích vô hướng qua ba cạnh xuất phát từ một đỉnh. a) Sai. b) Đúng. c) Sai. d) Đúng.}]{Cho hình lập phương $ABCD.A'B'C'D'$ có cạnh bằng $3$. Xét tính đúng sai của các mệnh đề sau:}{%
  \y{$\overrightarrow{DA'}\cdot \overrightarrow{A'B'}=-1$.}{s}
  \y{$|\overrightarrow{A'B}|^2=18$.}{d}
  \y{$|\overrightarrow{AD'}|^2=19$.}{s}
  \y{$|\overrightarrow{A'C}|^2=27$.}{d}}

\cauDS[id=C2-B6-D03-TH-051,muc=TH,loigiai={Tính các tích vô hướng qua ba cạnh xuất phát từ một đỉnh. a) Sai. b) Đúng. c) Đúng. d) Đúng.}]{Cho hình lập phương $EFGH.E'F'G'H'$ có cạnh bằng $2$. Xét tính đúng sai của các mệnh đề sau:}{%
  \y{$\overrightarrow{F'H'}\cdot \overrightarrow{E'H'}=3$.}{s}
  \y{$\overrightarrow{HE}\cdot \overrightarrow{HG}=0$.}{d}
  \y{$\overrightarrow{E'F}\cdot \overrightarrow{F'H}=0$.}{d}
  \y{$|\overrightarrow{FE}|^2=4$.}{d}}

\cauDS[id=C2-B6-D03-TH-052,muc=TH,loigiai={Tính các tích vô hướng qua ba cạnh xuất phát từ một đỉnh. a) Sai. b) Đúng. c) Sai. d) Đúng.}]{Cho tứ diện đều $SABC$ có cạnh bằng $1$. Xét tính đúng sai của các mệnh đề sau:}{%
  \y{$|\overrightarrow{SB}|^2=-1$.}{s}
  \y{$\overrightarrow{CS}\cdot \overrightarrow{SA}=-\dfrac{1}{2}$.}{d}
  \y{$|\overrightarrow{SC}|^2=-1$.}{s}
  \y{$\overrightarrow{BC}\cdot \overrightarrow{SA}=0$.}{d}}

\cauDS[id=C2-B6-D03-TH-053,muc=TH,loigiai={Tính các tích vô hướng qua ba cạnh xuất phát từ một đỉnh. a) Sai. b) Đúng. c) Sai. d) Sai.}]{Cho hình lập phương $ABCD.EFGH$ có cạnh bằng $2$. Xét tính đúng sai của các mệnh đề sau:}{%
  \y{Hai vectơ $\overrightarrow{DC}$ và $\overrightarrow{DF}$ vuông góc với nhau.}{s}
  \y{Hai vectơ $\overrightarrow{BF}$ và $\overrightarrow{DB}$ vuông góc với nhau.}{d}
  \y{$\overrightarrow{EG}\cdot \overrightarrow{HE}=-3$.}{s}
  \y{$|\overrightarrow{DH}|^2=-4$.}{s}}

\cauDS[id=C2-B6-D03-TH-054,muc=TH,loigiai={Tính các tích vô hướng qua ba cạnh xuất phát từ một đỉnh. a) Đúng. b) Sai. c) Sai. d) Sai.}]{Cho hình lập phương $ABCD.EFGH$ có cạnh bằng $1$. Xét tính đúng sai của các mệnh đề sau:}{%
  \y{$\overrightarrow{HE}\cdot \overrightarrow{CE}=1$.}{d}
  \y{$|\overrightarrow{CF}|^2=-2$.}{s}
  \y{$\overrightarrow{AD}\cdot \overrightarrow{FC}=-1$.}{s}
  \y{Hai vectơ $\overrightarrow{AD}$ và $\overrightarrow{BD}$ vuông góc với nhau.}{s}}

\cauDS[id=C2-B6-D03-TH-055,muc=TH,loigiai={Tính các tích vô hướng qua ba cạnh xuất phát từ một đỉnh. a) Đúng. b) Sai. c) Đúng. d) Đúng.}]{Cho hình hộp chữ nhật $EFGH.E'F'G'H'$ có $EF=3$, $EH=4$, $EE'=1$. Xét tính đúng sai của các mệnh đề sau:}{%
  \y{$|\overrightarrow{FG}|^2=16$.}{d}
  \y{$|\overrightarrow{H'G'}|^2=-9$.}{s}
  \y{$\overrightarrow{HE'}\cdot \overrightarrow{F'G}=-17$.}{d}
  \y{$\overrightarrow{HG}\cdot \overrightarrow{HE'}=0$.}{d}}

\cauDS[id=C2-B6-D03-TH-056,muc=TH,loigiai={Tính các tích vô hướng qua ba cạnh xuất phát từ một đỉnh. a) Đúng. b) Sai. c) Đúng. d) Sai.}]{Cho hình lập phương $ABCD.EFGH$ có cạnh bằng $3$. Xét tính đúng sai của các mệnh đề sau:}{%
  \y{$\overrightarrow{CB}\cdot \overrightarrow{GE}=9$.}{d}
  \y{$|\overrightarrow{BA}|^2=-9$.}{s}
  \y{$\overrightarrow{CF}\cdot \overrightarrow{AB}=0$.}{d}
  \y{$|\overrightarrow{CA}|^2=17$.}{s}}

\cauDS[id=C2-B6-D03-TH-057,muc=TH,loigiai={Tính các tích vô hướng qua ba cạnh xuất phát từ một đỉnh. a) Đúng. b) Sai. c) Đúng. d) Đúng.}]{Cho hình lập phương $EFGH.E'F'G'H'$ có cạnh bằng $2$. Xét tính đúng sai của các mệnh đề sau:}{%
  \y{$|\overrightarrow{GE}|^2=8$.}{d}
  \y{Hai vectơ $\overrightarrow{G'F'}$ và $\overrightarrow{HE}$ vuông góc với nhau.}{s}
  \y{$\overrightarrow{G'G}\cdot \overrightarrow{EF}=0$.}{d}
  \y{$|\overrightarrow{E'G'}|^2=8$.}{d}}

\cauDS[id=C2-B6-D03-TH-058,muc=TH,loigiai={Tính các tích vô hướng qua ba cạnh xuất phát từ một đỉnh. a) Sai. b) Đúng. c) Sai. d) Đúng.}]{Cho hình hộp chữ nhật $EFGH.E'F'G'H'$ có $EF=1$, $EH=1$, $EE'=3$. Xét tính đúng sai của các mệnh đề sau:}{%
  \y{$|\overrightarrow{FF'}|^2=8$.}{s}
  \y{$\overrightarrow{H'E}\cdot \overrightarrow{HF}=1$.}{d}
  \y{Hai vectơ $\overrightarrow{EH'}$ và $\overrightarrow{H'F'}$ vuông góc với nhau.}{s}
  \y{$\overrightarrow{E'E}\cdot \overrightarrow{EF'}=-9$.}{d}}

\cauDS[id=C2-B6-D03-TH-059,muc=TH,loigiai={Tính các tích vô hướng qua ba cạnh xuất phát từ một đỉnh. a) Đúng. b) Sai. c) Đúng. d) Đúng.}]{Cho hình lập phương $MNPQ.M'N'P'Q'$ có cạnh bằng $a$. Xét tính đúng sai của các mệnh đề sau:}{%
  \y{$\overrightarrow{PQ'}\cdot \overrightarrow{MP}=-a^2$.}{d}
  \y{Hai vectơ $\overrightarrow{PQ'}$ và $\overrightarrow{QP}$ vuông góc với nhau.}{s}
  \y{$|\overrightarrow{N'Q}|^2=3a^2$.}{d}
  \y{$|\overrightarrow{N'M'}|^2=a^2$.}{d}}

\cauDS[id=C2-B6-D03-TH-060,muc=TH,loigiai={Tính các tích vô hướng qua ba cạnh xuất phát từ một đỉnh. a) Đúng. b) Sai. c) Sai. d) Đúng.}]{Cho hình lập phương $ABCD.EFGH$ có cạnh bằng $a$. Xét tính đúng sai của các mệnh đề sau:}{%
  \y{$\overrightarrow{FD}\cdot \overrightarrow{BA}=a^2$.}{d}
  \y{$\overrightarrow{AF}\cdot \overrightarrow{DE}=0$.}{s}
  \y{$|\overrightarrow{BD}|^2=3a^2$.}{s}
  \y{Hai vectơ $\overrightarrow{EH}$ và $\overrightarrow{DH}$ vuông góc với nhau.}{d}}

\cauDS[id=C2-B6-D03-TH-061,muc=TH,loigiai={Tính các tích vô hướng qua ba cạnh xuất phát từ một đỉnh. a) Đúng. b) Sai. c) Sai. d) Đúng.}]{Cho hình lập phương $EFGH.E'F'G'H'$ có cạnh bằng $a$. Xét tính đúng sai của các mệnh đề sau:}{%
  \y{$\overrightarrow{HE'}\cdot \overrightarrow{HG'}=a^2$.}{d}
  \y{$\overrightarrow{HG'}\cdot \overrightarrow{FE}=a^2$.}{s}
  \y{$|\overrightarrow{E'H'}|^2=2a^2$.}{s}
  \y{$\overrightarrow{FG'}\cdot \overrightarrow{HG'}=a^2$.}{d}}

\cauDS[id=C2-B6-D03-TH-062,muc=TH,loigiai={Tính các tích vô hướng qua ba cạnh xuất phát từ một đỉnh. a) Đúng. b) Sai. c) Đúng. d) Sai.}]{Cho hình hộp chữ nhật $ABCD.A'B'C'D'$ có $AB=3$, $AD=4$, $AA'=3$. Xét tính đúng sai của các mệnh đề sau:}{%
  \y{$\overrightarrow{B'A'}\cdot \overrightarrow{A'C'}=-9$.}{d}
  \y{$\overrightarrow{D'B}\cdot \overrightarrow{DB}=26$.}{s}
  \y{$\overrightarrow{DB'}\cdot \overrightarrow{BB'}=9$.}{d}
  \y{$|\overrightarrow{CA}|^2=24$.}{s}}

\cauDS[id=C2-B6-D03-TH-063,muc=TH,loigiai={Tính các tích vô hướng qua ba cạnh xuất phát từ một đỉnh. a) Sai. b) Đúng. c) Sai. d) Sai.}]{Cho hình hộp chữ nhật $ABCD.EFGH$ có $AB=1$, $AD=1$, $AE=3$. Xét tính đúng sai của các mệnh đề sau:}{%
  \y{$\overrightarrow{BF}\cdot \overrightarrow{GB}=-8$.}{s}
  \y{$\overrightarrow{GC}\cdot \overrightarrow{AD}=0$.}{d}
  \y{Hai vectơ $\overrightarrow{AG}$ và $\overrightarrow{GH}$ vuông góc với nhau.}{s}
  \y{$|\overrightarrow{DE}|^2=9$.}{s}}

\cauDS[id=C2-B6-D03-TH-064,muc=TH,loigiai={Tính các tích vô hướng qua ba cạnh xuất phát từ một đỉnh. a) Đúng. b) Đúng. c) Đúng. d) Sai.}]{Cho hình lập phương $EFGH.E'F'G'H'$ có cạnh bằng $2$. Xét tính đúng sai của các mệnh đề sau:}{%
  \y{$|\overrightarrow{H'F}|^2=12$.}{d}
  \y{$|\overrightarrow{FF'}|^2=4$.}{d}
  \y{$\overrightarrow{H'H}\cdot \overrightarrow{GE'}=-4$.}{d}
  \y{$|\overrightarrow{EH'}|^2=9$.}{s}}

\cauTLN[id=C2-B6-D03-TH-065,muc=TH,loigiai={Phân tích theo ba cạnh xuất phát từ một đỉnh ta được tích vô hướng bằng $-1$.}]{Cho hình hộp chữ nhật $ABCD.A'B'C'D'$ có $AB=1$, $AD=4$, $AA'=3$. Tính $\overrightarrow{D'B}\cdot \overrightarrow{C'D'}$.}{-1}

\cauTLN[id=C2-B6-D03-TH-066,muc=TH,loigiai={Ta có $\overrightarrow{BA}+2\overrightarrow{CD}=-\vec{a}-2\vec{b}+2\vec{c}$ nên $|\overrightarrow{BA}+2\overrightarrow{CD}|^2=20$.}]{Cho tứ diện đều $ABCD$ có cạnh bằng $2$. Tính $|\overrightarrow{BA}+2\overrightarrow{CD}|^2$.}{20}

\cauTLN[id=C2-B6-D03-TH-067,muc=TH,loigiai={Phân tích theo ba cạnh xuất phát từ một đỉnh ta được tích vô hướng bằng $-1$.}]{Cho hình lập phương $ABCD.A'B'C'D'$ có cạnh bằng $1$. Tính $\overrightarrow{CD}\cdot \overrightarrow{A'B}$.}{-1}

\cauTLN[id=C2-B6-D03-TH-068,muc=TH,loigiai={Ta có $2\overrightarrow{FH}=-2\vec{a}+2\vec{b}$ nên $|2\overrightarrow{FH}|^2=20$.}]{Cho hình hộp chữ nhật $EFGH.E'F'G'H'$ có $EF=2$, $EH=1$, $EE'=1$. Tính $|2\overrightarrow{FH}|^2$.}{20}

\cauTLN[id=C2-B6-D03-TH-069,muc=TH,loigiai={Phân tích theo ba cạnh xuất phát từ một đỉnh ta được tích vô hướng bằng $4$.}]{Cho hình hộp chữ nhật $MNPQ.M'N'P'Q'$ có $MN=2$, $MQ=1$, $MM'=2$. Tính $\overrightarrow{MN}\cdot \overrightarrow{Q'N'}$.}{4}

\cauTLN[id=C2-B6-D03-TH-070,muc=TH,loigiai={Phân tích theo ba cạnh xuất phát từ một đỉnh ta được tích vô hướng bằng $0$.}]{Cho hình hộp chữ nhật $ABCD.EFGH$ có $AB=3$, $AD=2$, $AE=3$. Tính $\overrightarrow{FE}\cdot \overrightarrow{GF}$.}{0}

\cauTLN[id=C2-B6-D03-TH-071,muc=TH,loigiai={Ta có $2\overrightarrow{SC}-\overrightarrow{SA}=-\vec{a}+2\vec{c}$ nên $|2\overrightarrow{SC}-\overrightarrow{SA}|^2=12$.}]{Cho tứ diện đều $SABC$ có cạnh bằng $2$. Tính $|2\overrightarrow{SC}-\overrightarrow{SA}|^2$.}{12}

\cauTLN[id=C2-B6-D03-TH-072,muc=TH,loigiai={Ta có $-\overrightarrow{Q'N}=-\vec{a}+\vec{b}+\vec{c}$ nên $|-\overrightarrow{Q'N}|^2=3$.}]{Cho hình lập phương $MNPQ.M'N'P'Q'$ có cạnh bằng $1$. Tính $|-\overrightarrow{Q'N}|^2$.}{3}

\cauTLN[id=C2-B6-D03-TH-073,muc=TH,loigiai={Ta có $2\overrightarrow{BD'}=-2\vec{a}+2\vec{b}+2\vec{c}$ nên $|2\overrightarrow{BD'}|^2=48$.}]{Cho hình lập phương $ABCD.A'B'C'D'$ có cạnh bằng $2$. Tính $|2\overrightarrow{BD'}|^2$.}{48}

\cauTLN[id=C2-B6-D03-TH-074,muc=TH,loigiai={Ta có $-\overrightarrow{SC}+\overrightarrow{BA}=\vec{a}-\vec{b}-\vec{c}$ nên $|-\overrightarrow{SC}+\overrightarrow{BA}|^2=8$.}]{Cho tứ diện đều $SABC$ có cạnh bằng $2$. Tính $|-\overrightarrow{SC}+\overrightarrow{BA}|^2$.}{8}

\cauTLN[id=C2-B6-D03-TH-075,muc=TH,loigiai={Ta có $\overrightarrow{CB}-\overrightarrow{SC}=\vec{b}-2\vec{c}$ nên $|\overrightarrow{CB}-\overrightarrow{SC}|^2=3$.}]{Cho tứ diện đều $SABC$ có cạnh bằng $1$. Tính $|\overrightarrow{CB}-\overrightarrow{SC}|^2$.}{3}

\cauTLN[id=C2-B6-D03-TH-076,muc=TH,loigiai={Ta có $\overrightarrow{AS}+2\overrightarrow{AB}=-3\vec{a}+2\vec{b}$ nên $|\overrightarrow{AS}+2\overrightarrow{AB}|^2=7$.}]{Cho tứ diện đều $SABC$ có cạnh bằng $1$. Tính $|\overrightarrow{AS}+2\overrightarrow{AB}|^2$.}{7}

\cauTLN[id=C2-B6-D03-TH-077,muc=TH,loigiai={Ta có $-\overrightarrow{DB}+\overrightarrow{GC}=-\vec{a}+\vec{b}-\vec{c}$ nên $|-\overrightarrow{DB}+\overrightarrow{GC}|^2=3$.}]{Cho hình lập phương $ABCD.EFGH$ có cạnh bằng $1$. Tính $|-\overrightarrow{DB}+\overrightarrow{GC}|^2$.}{3}

\cauTLN[id=C2-B6-D03-TH-078,muc=TH,loigiai={Ta có $-\overrightarrow{B'D'}=\vec{a}-\vec{b}$ nên $|-\overrightarrow{B'D'}|^2=8$.}]{Cho hình lập phương $ABCD.A'B'C'D'$ có cạnh bằng $2$. Tính $|-\overrightarrow{B'D'}|^2$.}{8}

\cauTLN[id=C2-B6-D03-TH-079,muc=TH,loigiai={Ta có $\overrightarrow{GC}+\overrightarrow{DB}=\vec{a}-\vec{b}-\vec{c}$ nên $|\overrightarrow{GC}+\overrightarrow{DB}|^2=3$.}]{Cho hình lập phương $ABCD.EFGH$ có cạnh bằng $1$. Tính $|\overrightarrow{GC}+\overrightarrow{DB}|^2$.}{3}

\cauTLN[id=C2-B6-D03-TH-080,muc=TH,loigiai={Ta có $2\overrightarrow{BC}+\overrightarrow{BD}=-3\vec{a}+2\vec{b}+\vec{c}$ nên $|2\overrightarrow{BC}+\overrightarrow{BD}|^2=28$.}]{Cho tứ diện đều $ABCD$ có cạnh bằng $2$. Tính $|2\overrightarrow{BC}+\overrightarrow{BD}|^2$.}{28}

\cauTN[id=C2-B6-D03-VD-001,muc=VD,loigiai={Ta có $\overrightarrow{FD}\cdot \overrightarrow{EF}=-9$, $|\overrightarrow{FD}|^2=26$, $|\overrightarrow{EF}|^2=9$ nên $\cos(\overrightarrow{FD},\overrightarrow{EF})=-\dfrac{3\sqrt{26}}{26}$.}]{Cho hình hộp chữ nhật $ABCD.EFGH$ có $AB=3$, $AD=4$, $AE=1$. Giá trị $\cos(\overrightarrow{FD},\overrightarrow{EF})$ bằng}
  {$-\dfrac{3\sqrt{26}}{26}$.}
  {$-\dfrac{4\sqrt{26}}{39}$.}{$\dfrac{3\sqrt{26}}{26}$.}{$-\dfrac{5\sqrt{26}}{39}$.}

\cauTN[id=C2-B6-D03-VD-002,muc=VD,loigiai={Ta có $\overrightarrow{CB}\cdot \overrightarrow{CA}=2$, $|\overrightarrow{CB}|^2=4$, $|\overrightarrow{CA}|^2=4$ nên $\cos(\overrightarrow{CB},\overrightarrow{CA})=\dfrac{1}{2}$.}]{Cho tứ diện đều $ABCD$ có cạnh bằng $2$. Giá trị $\cos(\overrightarrow{CB},\overrightarrow{CA})$ bằng}
  {$\dfrac{1}{2}$.}
  {$\dfrac{3}{4}$.}{$-\dfrac{1}{2}$.}{$\dfrac{1}{4}$.}

\cauTN[id=C2-B6-D03-VD-003,muc=VD,loigiai={Ta có $\overrightarrow{HB}\cdot \overrightarrow{CE}=-21$, $|\overrightarrow{HB}|^2=29$, $|\overrightarrow{CE}|^2=29$ nên $\cos(\overrightarrow{HB},\overrightarrow{CE})=-\dfrac{21}{29}$.}]{Cho hình hộp chữ nhật $ABCD.EFGH$ có $AB=3$, $AD=2$, $AE=4$. Giá trị $\cos(\overrightarrow{HB},\overrightarrow{CE})$ bằng}
  {$-\dfrac{21}{29}$.}
  {$-\dfrac{20}{29}$.}{$\dfrac{21}{29}$.}{$-\dfrac{22}{29}$.}

\cauTN[id=C2-B6-D03-VD-004,muc=VD,loigiai={Ta có $\overrightarrow{NP}\cdot \overrightarrow{PM}=-\dfrac{9}{2}$, $|\overrightarrow{NP}|^2=9$, $|\overrightarrow{PM}|^2=9$ nên $\cos(\overrightarrow{NP},\overrightarrow{PM})=-\dfrac{1}{2}$.}]{Cho tứ diện đều $MNPQ$ có cạnh bằng $3$. Giá trị $\cos(\overrightarrow{NP},\overrightarrow{PM})$ bằng}
  {$-\dfrac{1}{2}$.}
  {$-\dfrac{7}{18}$.}{$\dfrac{1}{2}$.}{$-\dfrac{11}{18}$.}

\cauTN[id=C2-B6-D03-VD-005,muc=VD,loigiai={Ta có $\overrightarrow{Q'P}\cdot \overrightarrow{M'M}=9$, $|\overrightarrow{Q'P}|^2=13$, $|\overrightarrow{M'M}|^2=9$ nên $\cos(\overrightarrow{Q'P},\overrightarrow{M'M})=\dfrac{3\sqrt{13}}{13}$.}]{Cho hình hộp chữ nhật $MNPQ.M'N'P'Q'$ có $MN=2$, $MQ=4$, $MM'=3$. Giá trị $\cos(\overrightarrow{Q'P},\overrightarrow{M'M})$ bằng}
  {$\dfrac{3\sqrt{13}}{13}$.}
  {$\dfrac{10\sqrt{13}}{39}$.}{$-\dfrac{3\sqrt{13}}{13}$.}{$\dfrac{8\sqrt{13}}{39}$.}

\cauTN[id=C2-B6-D03-VD-006,muc=VD,loigiai={Ta có $\overrightarrow{AG}\cdot \overrightarrow{CA}=-2a^2$, $|\overrightarrow{AG}|^2=3a^2$, $|\overrightarrow{CA}|^2=2a^2$ nên $\cos(\overrightarrow{AG},\overrightarrow{CA})=-\dfrac{\sqrt{6}}{3}$.}]{Cho hình lập phương $ABCD.EFGH$ có cạnh bằng $a$. Giá trị $\cos(\overrightarrow{AG},\overrightarrow{CA})$ bằng}
  {$-\dfrac{\sqrt{6}}{3}$.}
  {$-\dfrac{\sqrt{6}}{6}$.}{$\dfrac{\sqrt{6}}{3}$.}{$0$.}

\cauTN[id=C2-B6-D03-VD-007,muc=VD,loigiai={Ta có $\overrightarrow{NP'}\cdot \overrightarrow{M'N}=-9$, $|\overrightarrow{NP'}|^2=18$, $|\overrightarrow{M'N}|^2=18$ nên $\cos(\overrightarrow{NP'},\overrightarrow{M'N})=-\dfrac{1}{2}$.}]{Cho hình lập phương $MNPQ.M'N'P'Q'$ có cạnh bằng $3$. Giá trị $\cos(\overrightarrow{NP'},\overrightarrow{M'N})$ bằng}
  {$-\dfrac{1}{2}$.}
  {$-\dfrac{4}{9}$.}{$\dfrac{1}{2}$.}{$-\dfrac{5}{9}$.}

\cauTN[id=C2-B6-D03-VD-008,muc=VD,loigiai={Ta có $\overrightarrow{E'E}\cdot \overrightarrow{GF'}=-4$, $|\overrightarrow{E'E}|^2=4$, $|\overrightarrow{GF'}|^2=8$ nên $\cos(\overrightarrow{E'E},\overrightarrow{GF'})=-\dfrac{\sqrt{2}}{2}$.}]{Cho hình lập phương $EFGH.E'F'G'H'$ có cạnh bằng $2$. Giá trị $\cos(\overrightarrow{E'E},\overrightarrow{GF'})$ bằng}
  {$-\dfrac{\sqrt{2}}{2}$.}
  {$-\dfrac{3\sqrt{2}}{8}$.}{$\dfrac{\sqrt{2}}{2}$.}{$-\dfrac{5\sqrt{2}}{8}$.}

\cauTN[id=C2-B6-D03-VD-009,muc=VD,loigiai={Ta có $\overrightarrow{A'D'}\cdot \overrightarrow{C'B}=-16$, $|\overrightarrow{A'D'}|^2=16$, $|\overrightarrow{C'B}|^2=17$ nên $\cos(\overrightarrow{A'D'},\overrightarrow{C'B})=-\dfrac{4\sqrt{17}}{17}$.}]{Cho hình hộp chữ nhật $ABCD.A'B'C'D'$ có $AB=4$, $AD=4$, $AA'=1$. Giá trị $\cos(\overrightarrow{A'D'},\overrightarrow{C'B})$ bằng}
  {$-\dfrac{4\sqrt{17}}{17}$.}
  {$-\dfrac{15\sqrt{17}}{68}$.}{$\dfrac{4\sqrt{17}}{17}$.}{$-\dfrac{7\sqrt{17}}{34}$.}

\cauTN[id=C2-B6-D03-VD-010,muc=VD,loigiai={Ta có $\overrightarrow{GD}\cdot \overrightarrow{HF}=-16$, $|\overrightarrow{GD}|^2=20$, $|\overrightarrow{HF}|^2=17$ nên $\cos(\overrightarrow{GD},\overrightarrow{HF})=-\dfrac{8\sqrt{85}}{85}$.}]{Cho hình hộp chữ nhật $ABCD.EFGH$ có $AB=4$, $AD=1$, $AE=2$. Giá trị $\cos(\overrightarrow{GD},\overrightarrow{HF})$ bằng}
  {$-\dfrac{8\sqrt{85}}{85}$.}
  {$-\dfrac{3\sqrt{85}}{34}$.}{$\dfrac{8\sqrt{85}}{85}$.}{$-\dfrac{\sqrt{85}}{10}$.}

\cauTN[id=C2-B6-D03-VD-011,muc=VD,loigiai={Ta có $\overrightarrow{BC}\cdot \overrightarrow{CB}=-9$, $|\overrightarrow{BC}|^2=9$, $|\overrightarrow{CB}|^2=9$ nên $\cos(\overrightarrow{BC},\overrightarrow{CB})=-1$.}]{Cho tứ diện đều $SABC$ có cạnh bằng $3$. Giá trị $\cos(\overrightarrow{BC},\overrightarrow{CB})$ bằng}
  {$-1$.}
  {$-\dfrac{8}{9}$.}{$1$.}{$-\dfrac{7}{9}$.}

\cauTN[id=C2-B6-D03-VD-012,muc=VD,loigiai={Ta có $\overrightarrow{GF}\cdot \overrightarrow{DB}=9$, $|\overrightarrow{GF}|^2=9$, $|\overrightarrow{DB}|^2=18$ nên $\cos(\overrightarrow{GF},\overrightarrow{DB})=\dfrac{\sqrt{2}}{2}$.}]{Cho hình lập phương $ABCD.EFGH$ có cạnh bằng $3$. Giá trị $\cos(\overrightarrow{GF},\overrightarrow{DB})$ bằng}
  {$\dfrac{\sqrt{2}}{2}$.}
  {$\dfrac{5\sqrt{2}}{9}$.}{$-\dfrac{\sqrt{2}}{2}$.}{$\dfrac{4\sqrt{2}}{9}$.}

\cauTN[id=C2-B6-D03-VD-013,muc=VD,loigiai={Ta có $\vec{a}\cdot\vec{b}=2\cdot1\cdot\cos 60^\circ=1$. Khai triển: $(-\vec{a}+\vec{b})\cdot(-\vec{a}+3\vec{b})=1\vec{a}^2-4\vec{a}\cdot\vec{b}+3\vec{b}^2=3$.}]{Cho hai vectơ $\vec{a}$, $\vec{b}$ có $|\vec{a}|=2$, $|\vec{b}|=1$ và $(\vec{a},\vec{b})=60^\circ$. Tích vô hướng $(-\vec{a}+\vec{b})\cdot(-\vec{a}+3\vec{b})$ bằng}
  {$3$.}
  {$2$.}{$4$.}{$1$.}

\cauTN[id=C2-B6-D03-VD-014,muc=VD,loigiai={Ta có $\vec{a}\cdot\vec{b}=3\cdot5\cdot\cos 90^\circ=0$. Khai triển: $(2\vec{a}+2\vec{b})\cdot(\vec{a}+\vec{b})=2\vec{a}^2+4\vec{a}\cdot\vec{b}+2\vec{b}^2=68$.}]{Cho hai vectơ $\vec{a}$, $\vec{b}$ có $|\vec{a}|=3$, $|\vec{b}|=5$ và $(\vec{a},\vec{b})=90^\circ$. Tích vô hướng $(2\vec{a}+2\vec{b})\cdot(\vec{a}+\vec{b})$ bằng}
  {$68$.}
  {$136$.}{$69$.}{$70$.}

\cauTN[id=C2-B6-D03-VD-015,muc=VD,loigiai={Ta có $\vec{a}\cdot\vec{b}=2\cdot3\cdot\cos 90^\circ=0$. Khai triển: $(\vec{a}-\vec{b})\cdot(2\vec{a}+3\vec{b})=2\vec{a}^2+1\vec{a}\cdot\vec{b}-3\vec{b}^2=-19$.}]{Cho hai vectơ $\vec{a}$, $\vec{b}$ có $|\vec{a}|=2$, $|\vec{b}|=3$ và $(\vec{a},\vec{b})=90^\circ$. Tích vô hướng $(\vec{a}-\vec{b})\cdot(2\vec{a}+3\vec{b})$ bằng}
  {$-19$.}
  {$-37$.}{$-38$.}{$-\dfrac{19}{2}$.}

\cauTN[id=C2-B6-D03-VD-016,muc=VD,loigiai={Ta có $\vec{a}\cdot\vec{b}=4\cdot1\cdot\cos 120^\circ=-2$. Khai triển: $(2\vec{a}+3\vec{b})\cdot(-\vec{a}+\vec{b})=-2\vec{a}^2-1\vec{a}\cdot\vec{b}+3\vec{b}^2=-27$.}]{Cho hai vectơ $\vec{a}$, $\vec{b}$ có $|\vec{a}|=4$, $|\vec{b}|=1$ và $(\vec{a},\vec{b})=120^\circ$. Tích vô hướng $(2\vec{a}+3\vec{b})\cdot(-\vec{a}+\vec{b})$ bằng}
  {$-27$.}
  {$-\dfrac{27}{2}$.}{$-54$.}{$-24$.}

\cauTN[id=C2-B6-D03-VD-017,muc=VD,loigiai={Ta có $\vec{a}\cdot\vec{b}=3\cdot2\cdot\cos 120^\circ=-3$. Khai triển: $(-\vec{a}+2\vec{b})\cdot(3\vec{a}+2\vec{b})=-3\vec{a}^2+4\vec{a}\cdot\vec{b}+4\vec{b}^2=-23$.}]{Cho hai vectơ $\vec{a}$, $\vec{b}$ có $|\vec{a}|=3$, $|\vec{b}|=2$ và $(\vec{a},\vec{b})=120^\circ$. Tích vô hướng $(-\vec{a}+2\vec{b})\cdot(3\vec{a}+2\vec{b})$ bằng}
  {$-23$.}
  {$-20$.}{$-22$.}{$-69$.}

\cauTN[id=C2-B6-D03-VD-018,muc=VD,loigiai={Ta có $\vec{a}\cdot\vec{b}=1\cdot3\cdot\cos 90^\circ=0$. Khai triển: $(\vec{a}+3\vec{b})\cdot(2\vec{a}+3\vec{b})=2\vec{a}^2+9\vec{a}\cdot\vec{b}+9\vec{b}^2=83$.}]{Cho hai vectơ $\vec{a}$, $\vec{b}$ có $|\vec{a}|=1$, $|\vec{b}|=3$ và $(\vec{a},\vec{b})=90^\circ$. Tích vô hướng $(\vec{a}+3\vec{b})\cdot(2\vec{a}+3\vec{b})$ bằng}
  {$83$.}
  {$82$.}{$167$.}{$85$.}

\cauTN[id=C2-B6-D03-VD-019,muc=VD,loigiai={Ta có $\vec{a}\cdot\vec{b}=1\cdot5\cdot\cos 90^\circ=0$. Khai triển: $(-\vec{a}-\vec{b})\cdot(3\vec{a}+3\vec{b})=-3\vec{a}^2-6\vec{a}\cdot\vec{b}-3\vec{b}^2=-78$.}]{Cho hai vectơ $\vec{a}$, $\vec{b}$ có $|\vec{a}|=1$, $|\vec{b}|=5$ và $(\vec{a},\vec{b})=90^\circ$. Tích vô hướng $(-\vec{a}-\vec{b})\cdot(3\vec{a}+3\vec{b})$ bằng}
  {$-78$.}
  {$-80$.}{$-234$.}{$-75$.}

\cauTN[id=C2-B6-D03-VD-020,muc=VD,loigiai={Ta có $\vec{a}\cdot\vec{b}=1\cdot1\cdot\cos 120^\circ=-\dfrac{1}{2}$. Khai triển: $(2\vec{a}+2\vec{b})\cdot(\vec{a}-\vec{b})=2\vec{a}^2+0\vec{a}\cdot\vec{b}-2\vec{b}^2=0$.}]{Cho hai vectơ $\vec{a}$, $\vec{b}$ có $|\vec{a}|=1$, $|\vec{b}|=1$ và $(\vec{a},\vec{b})=120^\circ$. Tích vô hướng $(2\vec{a}+2\vec{b})\cdot(\vec{a}-\vec{b})$ bằng}
  {$0$.}
  {$-2$.}{$-1$.}{$3$.}

\cauTN[id=C2-B6-D03-VD-021,muc=VD,loigiai={Ta có $\vec{a}\cdot\vec{b}=3\cdot3\cdot\cos 60^\circ=\dfrac{9}{2}$. Khai triển: $(\vec{a}+\vec{b})\cdot(2\vec{a}+2\vec{b})=2\vec{a}^2+4\vec{a}\cdot\vec{b}+2\vec{b}^2=54$.}]{Cho hai vectơ $\vec{a}$, $\vec{b}$ có $|\vec{a}|=3$, $|\vec{b}|=3$ và $(\vec{a},\vec{b})=60^\circ$. Tích vô hướng $(\vec{a}+\vec{b})\cdot(2\vec{a}+2\vec{b})$ bằng}
  {$54$.}
  {$57$.}{$108$.}{$-54$.}

\cauTN[id=C2-B6-D03-VD-022,muc=VD,loigiai={Ta có $\vec{a}\cdot\vec{b}=3\cdot2\cdot\cos 120^\circ=-3$. Khai triển: $(3\vec{a}+\vec{b})\cdot(\vec{a}+3\vec{b})=3\vec{a}^2+10\vec{a}\cdot\vec{b}+3\vec{b}^2=9$.}]{Cho hai vectơ $\vec{a}$, $\vec{b}$ có $|\vec{a}|=3$, $|\vec{b}|=2$ và $(\vec{a},\vec{b})=120^\circ$. Tích vô hướng $(3\vec{a}+\vec{b})\cdot(\vec{a}+3\vec{b})$ bằng}
  {$9$.}
  {$7$.}{$27$.}{$18$.}

\cauTN[id=C2-B6-D03-VD-023,muc=VD,loigiai={Ta có $\vec{a}\cdot\vec{b}=3\cdot1\cdot\cos 90^\circ=0$. Khai triển: $(3\vec{a}+3\vec{b})\cdot(-\vec{a}+2\vec{b})=-3\vec{a}^2+3\vec{a}\cdot\vec{b}+6\vec{b}^2=-21$.}]{Cho hai vectơ $\vec{a}$, $\vec{b}$ có $|\vec{a}|=3$, $|\vec{b}|=1$ và $(\vec{a},\vec{b})=90^\circ$. Tích vô hướng $(3\vec{a}+3\vec{b})\cdot(-\vec{a}+2\vec{b})$ bằng}
  {$-21$.}
  {$-23$.}{$-\dfrac{21}{2}$.}{$-20$.}

\cauTN[id=C2-B6-D03-VD-024,muc=VD,loigiai={Ta có $\vec{a}\cdot\vec{b}=4\cdot1\cdot\cos 90^\circ=0$. Khai triển: $(3\vec{a}-\vec{b})\cdot(\vec{a}+3\vec{b})=3\vec{a}^2+8\vec{a}\cdot\vec{b}-3\vec{b}^2=45$.}]{Cho hai vectơ $\vec{a}$, $\vec{b}$ có $|\vec{a}|=4$, $|\vec{b}|=1$ và $(\vec{a},\vec{b})=90^\circ$. Tích vô hướng $(3\vec{a}-\vec{b})\cdot(\vec{a}+3\vec{b})$ bằng}
  {$45$.}
  {$46$.}{$\dfrac{45}{2}$.}{$43$.}

\cauTN[id=C2-B6-D03-VD-025,muc=VD,loigiai={Biểu diễn $2\overrightarrow{P'M}+2\overrightarrow{MP}=-2\vec{c}$ theo ba cạnh xuất phát từ một đỉnh rồi dùng $|\vec{x}|^2=\vec{x}\cdot\vec{x}$ ta được $|2\overrightarrow{P'M}+2\overrightarrow{MP}|^2=36$, nên $|2\overrightarrow{P'M}+2\overrightarrow{MP}|=6$.}]{Cho hình lập phương $MNPQ.M'N'P'Q'$ có cạnh bằng $3$. Độ dài của vectơ $2\overrightarrow{P'M}+2\overrightarrow{MP}$ bằng}
  {$6$.}
  {$3\sqrt{3}$.}{$3\sqrt{6}$.}{$18$.}

\cauTN[id=C2-B6-D03-VD-026,muc=VD,loigiai={Biểu diễn $2\overrightarrow{CC'}-\overrightarrow{A'B}=-\vec{a}+3\vec{c}$ theo ba cạnh xuất phát từ một đỉnh rồi dùng $|\vec{x}|^2=\vec{x}\cdot\vec{x}$ ta được $|2\overrightarrow{CC'}-\overrightarrow{A'B}|^2=85$, nên $|2\overrightarrow{CC'}-\overrightarrow{A'B}|=\sqrt{85}$.}]{Cho hình hộp chữ nhật $ABCD.A'B'C'D'$ có $AB=2$, $AD=2$, $AA'=3$. Độ dài của vectơ $2\overrightarrow{CC'}-\overrightarrow{A'B}$ bằng}
  {$\sqrt{85}$.}
  {$\sqrt{13}$.}{$4\sqrt{10}$.}{$2\sqrt{10}$.}

\cauTN[id=C2-B6-D03-VD-027,muc=VD,loigiai={Biểu diễn $\overrightarrow{DC'}+2\overrightarrow{A'C}=3\vec{a}+2\vec{b}-\vec{c}$ theo ba cạnh xuất phát từ một đỉnh rồi dùng $|\vec{x}|^2=\vec{x}\cdot\vec{x}$ ta được $|\overrightarrow{DC'}+2\overrightarrow{A'C}|^2=169$, nên $|\overrightarrow{DC'}+2\overrightarrow{A'C}|=13$.}]{Cho hình hộp chữ nhật $ABCD.A'B'C'D'$ có $AB=4$, $AD=2$, $AA'=3$. Độ dài của vectơ $\overrightarrow{DC'}+2\overrightarrow{A'C}$ bằng}
  {$13$.}
  {$\sqrt{113}$.}{$4\sqrt{17}$.}{$2\sqrt{17}$.}

\cauTN[id=C2-B6-D03-VD-028,muc=VD,loigiai={Biểu diễn $\overrightarrow{D'B'}+2\overrightarrow{D'A}=\vec{a}-3\vec{b}-2\vec{c}$ theo ba cạnh xuất phát từ một đỉnh rồi dùng $|\vec{x}|^2=\vec{x}\cdot\vec{x}$ ta được $|\overrightarrow{D'B'}+2\overrightarrow{D'A}|^2=176$, nên $|\overrightarrow{D'B'}+2\overrightarrow{D'A}|=4\sqrt{11}$.}]{Cho hình hộp chữ nhật $ABCD.A'B'C'D'$ có $AB=4$, $AD=4$, $AA'=2$. Độ dài của vectơ $\overrightarrow{D'B'}+2\overrightarrow{D'A}$ bằng}
  {$4\sqrt{11}$.}
  {$4\sqrt{3}$.}{$2\sqrt{21}$.}{$4\sqrt{21}$.}

\cauTN[id=C2-B6-D03-VD-029,muc=VD,loigiai={Biểu diễn $-\overrightarrow{FA}-\overrightarrow{HB}=\vec{b}+2\vec{c}$ theo ba cạnh xuất phát từ một đỉnh rồi dùng $|\vec{x}|^2=\vec{x}\cdot\vec{x}$ ta được $|-\overrightarrow{FA}-\overrightarrow{HB}|^2=68$, nên $|-\overrightarrow{FA}-\overrightarrow{HB}|=2\sqrt{17}$.}]{Cho hình hộp chữ nhật $ABCD.EFGH$ có $AB=3$, $AD=2$, $AE=4$. Độ dài của vectơ $-\overrightarrow{FA}-\overrightarrow{HB}$ bằng}
  {$2\sqrt{17}$.}
  {$2\sqrt{10}$.}{$\sqrt{157}$.}{$13$.}

\cauTN[id=C2-B6-D03-VD-030,muc=VD,loigiai={Biểu diễn $\overrightarrow{AC}+\overrightarrow{AG}=2\vec{a}+2\vec{b}+\vec{c}$ theo ba cạnh xuất phát từ một đỉnh rồi dùng $|\vec{x}|^2=\vec{x}\cdot\vec{x}$ ta được $|\overrightarrow{AC}+\overrightarrow{AG}|^2=81$, nên $|\overrightarrow{AC}+\overrightarrow{AG}|=9$.}]{Cho hình hộp chữ nhật $ABCD.EFGH$ có $AB=3$, $AD=3$, $AE=3$. Độ dài của vectơ $\overrightarrow{AC}+\overrightarrow{AG}$ bằng}
  {$9$.}
  {$3\sqrt{19}$.}{$3$.}{$3\sqrt{22}$.}

\cauTN[id=C2-B6-D03-VD-031,muc=VD,loigiai={Biểu diễn $\overrightarrow{PN'}-\overrightarrow{QN'}=-\vec{a}$ theo ba cạnh xuất phát từ một đỉnh rồi dùng $|\vec{x}|^2=\vec{x}\cdot\vec{x}$ ta được $|\overrightarrow{PN'}-\overrightarrow{QN'}|^2=4$, nên $|\overrightarrow{PN'}-\overrightarrow{QN'}|=2$.}]{Cho hình hộp chữ nhật $MNPQ.M'N'P'Q'$ có $MN=2$, $MQ=3$, $MM'=1$. Độ dài của vectơ $\overrightarrow{PN'}-\overrightarrow{QN'}$ bằng}
  {$2$.}
  {$\sqrt{14}$.}{$2\sqrt{11}$.}{$\sqrt{26}$.}

\cauTN[id=C2-B6-D03-VD-032,muc=VD,loigiai={Biểu diễn $2\overrightarrow{N'M'}+\overrightarrow{M'N}=-\vec{a}-\vec{c}$ theo ba cạnh xuất phát từ một đỉnh rồi dùng $|\vec{x}|^2=\vec{x}\cdot\vec{x}$ ta được $|2\overrightarrow{N'M'}+\overrightarrow{M'N}|^2=2a^2$, nên $|2\overrightarrow{N'M'}+\overrightarrow{M'N}|=a\sqrt{2}$.}]{Cho hình lập phương $MNPQ.M'N'P'Q'$ có cạnh bằng $a$. Độ dài của vectơ $2\overrightarrow{N'M'}+\overrightarrow{M'N}$ bằng}
  {$a\sqrt{2}$.}
  {$a$.}{$a\sqrt{10}$.}{$2a$.}

\cauTN[id=C2-B6-D03-VD-033,muc=VD,loigiai={Biểu diễn $\overrightarrow{Q'N'}+2\overrightarrow{NM}=-\vec{a}-\vec{b}$ theo ba cạnh xuất phát từ một đỉnh rồi dùng $|\vec{x}|^2=\vec{x}\cdot\vec{x}$ ta được $|\overrightarrow{Q'N'}+2\overrightarrow{NM}|^2=25$, nên $|\overrightarrow{Q'N'}+2\overrightarrow{NM}|=5$.}]{Cho hình hộp chữ nhật $MNPQ.M'N'P'Q'$ có $MN=4$, $MQ=3$, $MM'=2$. Độ dài của vectơ $\overrightarrow{Q'N'}+2\overrightarrow{NM}$ bằng}
  {$5$.}
  {$3\sqrt{17}$.}{$6$.}{$3$.}

\cauTN[id=C2-B6-D03-VD-034,muc=VD,loigiai={Biểu diễn $\overrightarrow{Q'P'}+\overrightarrow{PM'}=-\vec{b}+\vec{c}$ theo ba cạnh xuất phát từ một đỉnh rồi dùng $|\vec{x}|^2=\vec{x}\cdot\vec{x}$ ta được $|\overrightarrow{Q'P'}+\overrightarrow{PM'}|^2=18$, nên $|\overrightarrow{Q'P'}+\overrightarrow{PM'}|=3\sqrt{2}$.}]{Cho hình lập phương $MNPQ.M'N'P'Q'$ có cạnh bằng $3$. Độ dài của vectơ $\overrightarrow{Q'P'}+\overrightarrow{PM'}$ bằng}
  {$3\sqrt{2}$.}
  {$3\sqrt{3}$.}{$3\sqrt{6}$.}{$9$.}

\cauTN[id=C2-B6-D03-VD-035,muc=VD,loigiai={Biểu diễn $\overrightarrow{B'D}+2\overrightarrow{AB}=\vec{a}+\vec{b}-\vec{c}$ theo ba cạnh xuất phát từ một đỉnh rồi dùng $|\vec{x}|^2=\vec{x}\cdot\vec{x}$ ta được $|\overrightarrow{B'D}+2\overrightarrow{AB}|^2=12$, nên $|\overrightarrow{B'D}+2\overrightarrow{AB}|=2\sqrt{3}$.}]{Cho hình lập phương $ABCD.A'B'C'D'$ có cạnh bằng $2$. Độ dài của vectơ $\overrightarrow{B'D}+2\overrightarrow{AB}$ bằng}
  {$2\sqrt{3}$.}
  {$4\sqrt{2}$.}{$2\sqrt{11}$.}{$2\sqrt{2}$.}

\cauTN[id=C2-B6-D03-VD-036,muc=VD,loigiai={Biểu diễn $\overrightarrow{CA'}+\overrightarrow{C'A'}=-2\vec{a}-2\vec{b}+\vec{c}$ theo ba cạnh xuất phát từ một đỉnh rồi dùng $|\vec{x}|^2=\vec{x}\cdot\vec{x}$ ta được $|\overrightarrow{CA'}+\overrightarrow{C'A'}|^2=56$, nên $|\overrightarrow{CA'}+\overrightarrow{C'A'}|=2\sqrt{14}$.}]{Cho hình hộp chữ nhật $ABCD.A'B'C'D'$ có $AB=3$, $AD=1$, $AA'=4$. Độ dài của vectơ $\overrightarrow{CA'}+\overrightarrow{C'A'}$ bằng}
  {$2\sqrt{14}$.}
  {$\sqrt{154}$.}{$4$.}{$\sqrt{106}$.}

\cauDS[id=C2-B6-D03-VD-037,muc=VD,loigiai={Tính các tích vô hướng qua ba cạnh xuất phát từ một đỉnh. a) Sai. b) Đúng. c) Sai. d) Đúng.}]{Cho hình lập phương $ABCD.EFGH$ có cạnh bằng $1$. Xét tính đúng sai của các mệnh đề sau:}{%
  \y{$|\overrightarrow{GE}|^2=1$.}{s}
  \y{$\overrightarrow{CA}\cdot \overrightarrow{AB}=-1$.}{d}
  \y{Hai vectơ $\overrightarrow{AF}$ và $\overrightarrow{DB}$ vuông góc với nhau.}{s}
  \y{Hai vectơ $\overrightarrow{BE}$ và $\overrightarrow{AF}$ vuông góc với nhau.}{d}}

\cauDS[id=C2-B6-D03-VD-038,muc=VD,loigiai={Tính các tích vô hướng qua ba cạnh xuất phát từ một đỉnh. a) Sai. b) Đúng. c) Sai. d) Đúng.}]{Cho hình hộp chữ nhật $ABCD.EFGH$ có $AB=4$, $AD=2$, $AE=1$. Xét tính đúng sai của các mệnh đề sau:}{%
  \y{Hai vectơ $\overrightarrow{HC}$ và $\overrightarrow{AC}$ vuông góc với nhau.}{s}
  \y{$\overrightarrow{DC}\cdot \overrightarrow{EF}=16$.}{d}
  \y{$\overrightarrow{CE}\cdot \overrightarrow{EA}=0$.}{s}
  \y{$|\overrightarrow{EC}|^2=21$.}{d}}

\cauDS[id=C2-B6-D03-VD-039,muc=VD,loigiai={Tính các tích vô hướng qua ba cạnh xuất phát từ một đỉnh. a) Sai. b) Sai. c) Sai. d) Đúng.}]{Cho hình lập phương $EFGH.E'F'G'H'$ có cạnh bằng $2$. Xét tính đúng sai của các mệnh đề sau:}{%
  \y{$|\overrightarrow{G'F'}|^2=5$.}{s}
  \y{$\overrightarrow{F'G}\cdot \overrightarrow{EG}=-4$.}{s}
  \y{$|\overrightarrow{GF'}|^2=-8$.}{s}
  \y{Hai vectơ $\overrightarrow{E'G}$ và $\overrightarrow{H'F'}$ vuông góc với nhau.}{d}}

\cauDS[id=C2-B6-D03-VD-040,muc=VD,loigiai={Tính các tích vô hướng qua ba cạnh xuất phát từ một đỉnh. a) Sai. b) Đúng. c) Đúng. d) Đúng.}]{Cho tứ diện đều $ABCD$ có cạnh bằng $2$. Xét tính đúng sai của các mệnh đề sau:}{%
  \y{$\overrightarrow{AB}\cdot \overrightarrow{CD}=-1$.}{s}
  \y{$\overrightarrow{CB}\cdot \overrightarrow{DB}=2$.}{d}
  \y{$|\overrightarrow{CB}|^2=4$.}{d}
  \y{$|\overrightarrow{BD}|^2=4$.}{d}}

\cauDS[id=C2-B6-D03-VD-041,muc=VD,loigiai={Tính các tích vô hướng qua ba cạnh xuất phát từ một đỉnh. a) Đúng. b) Sai. c) Đúng. d) Sai.}]{Cho tứ diện đều $SABC$ có cạnh bằng $2$. Xét tính đúng sai của các mệnh đề sau:}{%
  \y{$\overrightarrow{BC}\cdot \overrightarrow{BS}=2$.}{d}
  \y{$|\overrightarrow{AS}|^2=5$.}{s}
  \y{Hai vectơ $\overrightarrow{BA}$ và $\overrightarrow{SC}$ vuông góc với nhau.}{d}
  \y{$|\overrightarrow{SA}|^2=5$.}{s}}

\cauDS[id=C2-B6-D03-VD-042,muc=VD,loigiai={Tính các tích vô hướng qua ba cạnh xuất phát từ một đỉnh. a) Sai. b) Đúng. c) Sai. d) Sai.}]{Cho tứ diện đều $SABC$ có cạnh bằng $2$. Xét tính đúng sai của các mệnh đề sau:}{%
  \y{Hai vectơ $\overrightarrow{SA}$ và $\overrightarrow{BA}$ vuông góc với nhau.}{s}
  \y{$\overrightarrow{AC}\cdot \overrightarrow{AB}=2$.}{d}
  \y{Hai vectơ $\overrightarrow{CA}$ và $\overrightarrow{CB}$ vuông góc với nhau.}{s}
  \y{$|\overrightarrow{SB}|^2=3$.}{s}}

\cauDS[id=C2-B6-D03-VD-043,muc=VD,loigiai={Tính các tích vô hướng qua ba cạnh xuất phát từ một đỉnh. a) Sai. b) Đúng. c) Đúng. d) Đúng.}]{Cho tứ diện đều $MNPQ$ có cạnh bằng $3$. Xét tính đúng sai của các mệnh đề sau:}{%
  \y{$|\overrightarrow{NM}|^2=-9$.}{s}
  \y{$|\overrightarrow{PN}|^2=9$.}{d}
  \y{$\overrightarrow{PQ}\cdot \overrightarrow{QM}=-\dfrac{9}{2}$.}{d}
  \y{$|\overrightarrow{PQ}|^2=9$.}{d}}

\cauDS[id=C2-B6-D03-VD-044,muc=VD,loigiai={Tính các tích vô hướng qua ba cạnh xuất phát từ một đỉnh. a) Sai. b) Sai. c) Đúng. d) Sai.}]{Cho hình lập phương $ABCD.EFGH$ có cạnh bằng $a$. Xét tính đúng sai của các mệnh đề sau:}{%
  \y{$|\overrightarrow{EB}|^2=3a^2$.}{s}
  \y{$|\overrightarrow{BG}|^2=a^2$.}{s}
  \y{$|\overrightarrow{EB}|^2=2a^2$.}{d}
  \y{Hai vectơ $\overrightarrow{EC}$ và $\overrightarrow{EB}$ vuông góc với nhau.}{s}}

\cauDS[id=C2-B6-D03-VD-045,muc=VD,loigiai={Tính các tích vô hướng qua ba cạnh xuất phát từ một đỉnh. a) Sai. b) Đúng. c) Sai. d) Sai.}]{Cho hình hộp chữ nhật $MNPQ.M'N'P'Q'$ có $MN=2$, $MQ=3$, $MM'=4$. Xét tính đúng sai của các mệnh đề sau:}{%
  \y{$|\overrightarrow{MN'}|^2=-20$.}{s}
  \y{$|\overrightarrow{N'N}|^2=16$.}{d}
  \y{$|\overrightarrow{N'M}|^2=19$.}{s}
  \y{$|\overrightarrow{QM'}|^2=26$.}{s}}

\cauDS[id=C2-B6-D03-VD-046,muc=VD,loigiai={Tính các tích vô hướng qua ba cạnh xuất phát từ một đỉnh. a) Đúng. b) Đúng. c) Sai. d) Đúng.}]{Cho hình lập phương $MNPQ.M'N'P'Q'$ có cạnh bằng $3$. Xét tính đúng sai của các mệnh đề sau:}{%
  \y{$|\overrightarrow{P'Q}|^2=18$.}{d}
  \y{$|\overrightarrow{P'Q'}|^2=9$.}{d}
  \y{$|\overrightarrow{QM}|^2=8$.}{s}
  \y{$|\overrightarrow{Q'P'}|^2=9$.}{d}}

\cauDS[id=C2-B6-D03-VD-047,muc=VD,loigiai={Tính các tích vô hướng qua ba cạnh xuất phát từ một đỉnh. a) Đúng. b) Sai. c) Sai. d) Đúng.}]{Cho hình lập phương $EFGH.E'F'G'H'$ có cạnh bằng $1$. Xét tính đúng sai của các mệnh đề sau:}{%
  \y{$\overrightarrow{FE'}\cdot \overrightarrow{F'G}=-1$.}{d}
  \y{Hai vectơ $\overrightarrow{HF}$ và $\overrightarrow{EH'}$ vuông góc với nhau.}{s}
  \y{$|\overrightarrow{HG'}|^2=-2$.}{s}
  \y{$|\overrightarrow{H'H}|^2=1$.}{d}}

\cauDS[id=C2-B6-D03-VD-048,muc=VD,loigiai={Tính các tích vô hướng qua ba cạnh xuất phát từ một đỉnh. a) Đúng. b) Sai. c) Đúng. d) Đúng.}]{Cho tứ diện đều $ABCD$ có cạnh bằng $2$. Xét tính đúng sai của các mệnh đề sau:}{%
  \y{$\overrightarrow{DB}\cdot \overrightarrow{AD}=-2$.}{d}
  \y{Hai vectơ $\overrightarrow{BA}$ và $\overrightarrow{AD}$ vuông góc với nhau.}{s}
  \y{$|\overrightarrow{CA}|^2=4$.}{d}
  \y{Hai vectơ $\overrightarrow{AD}$ và $\overrightarrow{BC}$ vuông góc với nhau.}{d}}

\cauTLN[id=C2-B6-D03-VD-049,muc=VD,loigiai={Phân tích theo ba cạnh xuất phát từ một đỉnh ta được tích vô hướng bằng $0$.}]{Cho tứ diện đều $MNPQ$ có cạnh bằng $1$. Tính $\overrightarrow{PN}\cdot \overrightarrow{MQ}$.}{0}

\cauTLN[id=C2-B6-D03-VD-050,muc=VD,loigiai={Phân tích theo ba cạnh xuất phát từ một đỉnh ta được tích vô hướng bằng $-2$.}]{Cho tứ diện đều $SABC$ có cạnh bằng $2$. Tính $\overrightarrow{CB}\cdot \overrightarrow{AC}$.}{-2}

\cauTLN[id=C2-B6-D03-VD-051,muc=VD,loigiai={Phân tích theo ba cạnh xuất phát từ một đỉnh ta được tích vô hướng bằng $-9$.}]{Cho hình lập phương $ABCD.EFGH$ có cạnh bằng $3$. Tính $\overrightarrow{FE}\cdot \overrightarrow{HB}$.}{-9}

\cauTLN[id=C2-B6-D03-VD-052,muc=VD,loigiai={Ta có $\overrightarrow{P'P}+2\overrightarrow{P'N'}=-2\vec{b}-\vec{c}$ nên $|\overrightarrow{P'P}+2\overrightarrow{P'N'}|^2=37$.}]{Cho hình hộp chữ nhật $MNPQ.M'N'P'Q'$ có $MN=3$, $MQ=3$, $MM'=1$. Tính $|\overrightarrow{P'P}+2\overrightarrow{P'N'}|^2$.}{37}

\cauTLN[id=C2-B6-D03-VD-053,muc=VD,loigiai={Ta có $\overrightarrow{AD}=\vec{c}$ nên $|\overrightarrow{AD}|^2=4$.}]{Cho tứ diện đều $ABCD$ có cạnh bằng $2$. Tính $|\overrightarrow{AD}|^2$.}{4}

\cauTLN[id=C2-B6-D03-VD-054,muc=VD,loigiai={Phân tích theo ba cạnh xuất phát từ một đỉnh ta được tích vô hướng bằng $0$.}]{Cho hình lập phương $ABCD.A'B'C'D'$ có cạnh bằng $2$. Tính $\overrightarrow{C'D'}\cdot \overrightarrow{D'A'}$.}{0}

\cauTLN[id=C2-B6-D03-VD-055,muc=VD,loigiai={Ta có $2\overrightarrow{CB}=-2\vec{b}$ nên $|2\overrightarrow{CB}|^2=36$.}]{Cho hình hộp chữ nhật $ABCD.EFGH$ có $AB=4$, $AD=3$, $AE=2$. Tính $|2\overrightarrow{CB}|^2$.}{36}

\cauTLN[id=C2-B6-D03-VD-056,muc=VD,loigiai={Phân tích theo ba cạnh xuất phát từ một đỉnh ta được tích vô hướng bằng $-1$.}]{Cho hình lập phương $ABCD.EFGH$ có cạnh bằng $1$. Tính $\overrightarrow{ED}\cdot \overrightarrow{CH}$.}{-1}

\cauTLN[id=C2-B6-D03-VD-057,muc=VD,loigiai={Phân tích theo ba cạnh xuất phát từ một đỉnh ta được tích vô hướng bằng $-1$.}]{Cho hình lập phương $EFGH.E'F'G'H'$ có cạnh bằng $1$. Tính $\overrightarrow{G'E}\cdot \overrightarrow{FH'}$.}{-1}

\cauTLN[id=C2-B6-D03-VD-058,muc=VD,loigiai={Ta có $2\overrightarrow{D'C}=2\vec{a}-2\vec{c}$ nên $|2\overrightarrow{D'C}|^2=72$.}]{Cho hình lập phương $ABCD.A'B'C'D'$ có cạnh bằng $3$. Tính $|2\overrightarrow{D'C}|^2$.}{72}

\cauTLN[id=C2-B6-D03-VD-059,muc=VD,loigiai={Ta có $\overrightarrow{Q'N}=\vec{a}-\vec{b}-\vec{c}$ nên $|\overrightarrow{Q'N}|^2=12$.}]{Cho hình lập phương $MNPQ.M'N'P'Q'$ có cạnh bằng $2$. Tính $|\overrightarrow{Q'N}|^2$.}{12}

\cauTLN[id=C2-B6-D03-VD-060,muc=VD,loigiai={Phân tích theo ba cạnh xuất phát từ một đỉnh ta được tích vô hướng bằng $0$.}]{Cho hình hộp chữ nhật $ABCD.A'B'C'D'$ có $AB=3$, $AD=2$, $AA'=3$. Tính $\overrightarrow{BA'}\cdot \overrightarrow{AC'}$.}{0}
